\documentclass[output=paper]{langscibook}
\author{Authorname\orcid{}\affiliation{}}
\title{Title}
\abstract{Abstract}
\IfFileExists{../localcommands.tex}{
  \addbibresource{../localbibliography.bib}
  \usepackage{tabularx,multicol}
\usepackage[hyphens]{url}
\urlstyle{same}

\usepackage{langsci-optional}
\usepackage{langsci-lgr}
\usepackage{langsci-gb4e}

\usepackage{soul}
% \usepackage{booktabs}
% \usepackage{geometry}
\usepackage{multirow}
% \usepackage{makecell} %safely breaks a line inside a table cell
% \usepackage{tablefootnote} %footnotes in table captions
\usepackage{enumerate}
\usepackage{tikz}
\usepackage{array}
\usepackage[shortlabels]{enumitem}

%for having several figures on one line, turning words in 90°angles
% \usepackage{graphicx}
\usepackage{subcaption}

% \usepackage{caption} %for making the caption line longer in concrete cases

% \usepackage{rotating} %rotate a whole page; for large tables

\usepackage{fontspec}
\newfontfamily\charis{Charis SIL} %for being able to use the symbol ‿ and a better command of some of the boldfaced diacritics in 01
\newfontfamily\japanesefont{Noto Serif CJK JP} % for Japanese script
% \usepackage{tipa}

% \usepackage{needspace} %to keep exs on one page

\usepackage{xcolor}

\usepackage{pgfplots}
\pgfplotsset{compat=1.18}

%%%%%   AVMs for 02	%%%%%
\usepackage{langsci-avm}
\avmsetup{switch = ?}


\usepackage{langsci-branding}

  %\newcommand*{\orcid}{}


\makeatletter
\let\thetitle\@title
\let\theauthor\@author
\makeatother

\newcommand{\togglepaper}[1][0]{
   \bibliography{../localbibliography}
   \papernote{\scriptsize\normalfont
     \theauthor.
     \titleTemp.
     To appear in:
     E. Di Tor \& Herr Rausgeberin (ed.).
     Booktitle in localcommands.tex.
     Berlin: Language Science Press. [preliminary page numbering]
   }
   \pagenumbering{roman}
   \setcounter{chapter}{#1}
   \addtocounter{chapter}{-1}
}

\newbool{bookcompile}
\booltrue{bookcompile}
\newcommand{\bookorchapter}[2]{\ifbool{bookcompile}{#1}{#2}}



\renewcommand{\thempfootnote}{\fnsymbol{mpfootnote}} % use * for footnotes in float env

% \newcolumntype{L}[1]{>{\raggedright\arraybackslash}p{#1}} %left-aligned (non-justified) text in tables
% \newcolumntype{R}[1]{>{\raggedleft\arraybackslash}p{#1}} % right-aligned (non-justified) text in tables

\definecolor{customgreen}{RGB}{27,158,119}
\definecolor{custompink}{RGB}{231,42,138}
\definecolor{customblue}{RGB}{117,112,179}
\definecolor{customlightblue}{RGB}{143,170,220}
\definecolor{customdarkblue}{RGB}{68,114,196}
\definecolor{customyellow}{RGB}{225,192,0}
\definecolor{customorange}{RGB}{233,137,21}
\definecolor{customgray}{RGB}{229,229,229}


% Left-aligned indexes
% % \makeatletter
% % \patchcmd{\theindex}{\twocolumn}{\raggedright\twocolumn}{}{}
% % \makeatother

%to make Japanese scripts in the bibliography smaller
\newcommand{\japbib}[1]{{\japanesefont\small #1}}

%%%%%%%%%%%%%%%%%%%	MACROS for 02	%%%%%%%%%%%%%%%%

%%% Formatting	%%%%%

\newcommand{\textex}[1]{\textit{#1}} 		 %for formatting linguistic examples in-text%
%command-shift X

\newcommand{\lxm}[1]{\textsc{#1}}  		%for formatting names of lexemes
%command-option-shift L

%%%%%%%%% Abbreviations	%%%%%%

\newcommand{\featname}[1]{\textsc{#1}}
\newcommand{\featspec}[2]{\featname{#1}:{#2}}
\newcommand{\feat}[2]{\textsc{#1}:{#2}}  				% \feat{FEATURE}{value}  sc's
\newcommand{\featnameonly}[1]{\textsc{#1}} 			% \featnameonly{featurename}  sc's

\newcommand{\pr}{$^{\prime}$}

\newcommand{\Corr}{\textit{Corr}}
\newcommand{\propm}{\textit{pm}}

\newcommand{\PC}{Π$_{C}$}

\newcommand{\PF}{Π$_{F}$}
\newcommand{\PR}{Π$_{R}$}

\newcommand{\lex}{$\mathcal{L}$}
 
  %% hyphenation points for line breaks
%% Normally, automatic hyphenation in LaTeX is very good
%% If a word is mis-hyphenated, add it to this file
%%
%% add information to TeX file before \begin{document} with:
%% %% hyphenation points for line breaks
%% Normally, automatic hyphenation in LaTeX is very good
%% If a word is mis-hyphenated, add it to this file
%%
%% add information to TeX file before \begin{document} with:
%% %% hyphenation points for line breaks
%% Normally, automatic hyphenation in LaTeX is very good
%% If a word is mis-hyphenated, add it to this file
%%
%% add information to TeX file before \begin{document} with:
%% \include{localhyphenation}
\hyphenation{
    Al-ex-an-der
    Ber-ber
    chan-ges
    cir-cum-stan-ces
    coin-age
    Fraj-zyngier
    ita-liano
    Ja-pa-nese
    Kö-nig
    Krzy-żanowski
    li-te-ra-tur-no-go
    mi-zen-kei
    par-a-digm
    phe-nom-e-non
    ren-tai-shi
    Slo-vene
    Spen-cer
    Stoj-kov
    Ver-bi-ca-re-se
    wide-spread
}

\hyphenation{
    Al-ex-an-der
    Ber-ber
    chan-ges
    cir-cum-stan-ces
    coin-age
    Fraj-zyngier
    ita-liano
    Ja-pa-nese
    Kö-nig
    Krzy-żanowski
    li-te-ra-tur-no-go
    mi-zen-kei
    par-a-digm
    phe-nom-e-non
    ren-tai-shi
    Slo-vene
    Spen-cer
    Stoj-kov
    Ver-bi-ca-re-se
    wide-spread
}

\hyphenation{
    Al-ex-an-der
    Ber-ber
    chan-ges
    cir-cum-stan-ces
    coin-age
    Fraj-zyngier
    ita-liano
    Ja-pa-nese
    Kö-nig
    Krzy-żanowski
    li-te-ra-tur-no-go
    mi-zen-kei
    par-a-digm
    phe-nom-e-non
    ren-tai-shi
    Slo-vene
    Spen-cer
    Stoj-kov
    Ver-bi-ca-re-se
    wide-spread
}
 
  \togglepaper[1]%%chapternumber
}{}

\begin{document}
\maketitle 
%\shorttitlerunninghead{}%%use this for an abridged title in the page headers


\textbf{The} \textbf{dog} \textbf{didn’t} \textbf{bark,} \textbf{the} \textbf{noun} \textbf{didn’t} \textbf{inflect:}

\textbf{a} \textbf{typology} \textbf{of} \textbf{significant} \textbf{absences}\footnote{\textrm{A version was read at the Workshop “Uninflectedness” DGfS Cologne, \citealt{March2023}. The support of the ESRC (UK), grant: ES/R00837X/1 ‘Optimal categorisation: the origin and nature of gender from a psycholinguistic perspective’, and of the Leverhulme Trust (UK), grant RPG-2020-078 ‘Declining Case: Inflectional Loss in Progress’, is gratefully acknowledged. At various times the paper has benefitted from constructive comments by Jenny Audring, Matthew Baerman, Oliver Bond, Wayles Brown, Patricia Cabredo-Hofherr, Marina Chumakina, Sebastian Fedden, John Hutchinson, Steven Kaye, Yuni Kim, Masha Kyuseva, Ljubomir Popović, Erich Round, Andrew Spencer, Nadia Stronks, Greg Stump, Anna Thornton and Jurriaan Wiegertjes. I thank the referees for the improvements they suggested, and Penny Everson and Lisa Mack for their help in the preparation of the manuscript. My ORCID is 0000-0002-2667-9870.}}

Greville G. Corbett

Surrey Morphology Group \& Department of Human Behavior, Ecology and Culture, Max Planck Institute for Evolutionary Anthropology, Leipzig


\textit{Uninflectability is notable because it defies expectations and holds theoretical significance. The author questions why inflection is typically expected and why its absence is important. Uninflectability is distinguished from related phenomena like syncretism and defectiveness. Rather than viewing uninflectability as absolute, the text argues for a scale between fully inflected and uninflected items, where parts of a paradigm may remain uninflected. Using Canonical Typology, the author outlines four criteria to investigate uninflectability: \REF{ex:key:1} the extent of uninflectability across items, \REF{ex:key:2} within individual items, where paradigms may be partially uninflected, \REF{ex:key:3} external conditions influencing uninflectability, and \REF{ex:key:4} consequences imposed by uninflectable items. The analysis draws on Slavonic and Nakh-Dagestanian languages to refine definitions and explore the range of uninflectability, emphasizing the complexity and typological importance of uninflectable items.}


\section{\textbf{Introduction:} \textbf{Unexpected} \textbf{absence}}

Inspector Gregory of Scotland Yard calls in Sherlock Holmes to advise on the disappearance of a racehorse:\footnote{In \textit{Silver Blaze} from \textit{The Memoirs of Sherlock Holmes}, Arthur Conan Doyle, 1894.}  

“Is there any point to which you would wish to draw my attention?” \\
“To the curious incident of the dog in the night-time.” \\
“The dog did nothing in the night-time.” \\
“That was the curious incident,” remarked Sherlock Holmes.

There was an absence of the expected behaviour (barking), which Holmes points out as significant, and indeed leads to his solving the case. The interest of uninflectability is that our expectations are not met, and that this significant absence may reveal more than we could learn from an expected presence. 

\subsection{\textbf{Expectations}}

In inflectional morphology, as in linguistics generally, we have names for the unexpected. We have ‘suppletion’, ‘defectiveness’, and so on, but we have no special term for items that fully meet our expectations (‘regular’ comes closest, but even here items may be regular in one respect and irregular in another). For uninflectability, it is the ‘inflectional expectation’ \citep[146]{Spencer2020} that makes it nameworthy.

\subsection{\textbf{Intellectual} \textbf{housekeeping}}

Our expectations can be a starting point. To go further I look for clear definitions (and the attendant terms). I use ‘uninflected’ for items which do not bear inflection. This is the superordinate term; within ‘uninflected’, \citet[142]{Spencer2020} makes a helpful distinction between ‘uninflecting’ and ‘uninflectable’. Words that belong to a class which, in the given language, are not associated with a paradigm of inflection are ‘uninflecting’. (We do not expect them to inflect, any more than we expect frogs to bark.) More interesting are words which belong to a class where we find inflection in a given language, yet they fail to inflect; these are termed ‘uninflectable’. (We do expect dogs to bark, and when the dog in \textit{Silver Blaze} fails to bark, this is noteworthy.) We can specify these distinctions exactly, giving a logical end point to calibrate from, to yield a canonical typology (\sectref{sec:key:2}). I also discuss phenomena related to uninflectedness, since it is useful both to attempt to draw boundaries, and to recognize the interest of phenomena which resist classification.

\subsection{\textbf{Key} \textbf{data}}

Our data are taken from a range of sources, with two families figuring prominently. First, the Slavonic languages provide several examples. These languages are promising in having complex inflectional morphology, and a research tradition on uninflectability.\footnote{ \textrm{Note, for instance, Doleschal (2000-2001, 2002a, b, 2008, 2025). Uninflectedness figures relatively little in large surveys of the family, though the Comrie \& Corbett volume (1993) includes several relevant points in descriptions of particular languages, and Sussex \& \citet[250-251]{Cubberley2006} have brief, helpful notes. From a rather different perspective, \citet{Maldjieva1995} considers how the uninflecting items may be assigned a part of speech automatically, given that morphological indicators are lacking.}} The languages share many characteristics (which is valuable since it eliminates some variables) and yet they show interesting differences in terms of uninflectability. For instance, in all of them we expect nouns to inflect (for number, and in most languages for case too), but we also find instances of uninflectable nouns. As a first instance, consider the data in \tabref{tab:key:1}.


\begin{tabularx}{\textwidth}{XXXXXX}

\lsptoprule
%\hhline%%replace by cmidrule{~-----}
\multicolumn{1}{c}{} & Russian & Ukrainian & Polish & Czech & Serbo-Croat\\
\textsc{singular} & kakadu & kakadu & kakadu & kakadu & kakadu\\
\textsc{plural} & kakadu & kakadu & kakadu & kakadu-ové & kakadu-i\\
\lspbottomrule
\end{tabularx}
%%please move \begin{table} just above \begin{tabular
\begin{table}
\caption{\textit{Kakadu} ‘cockatoo’ in a selection of Slavonic languages}
\label{tab:key:1}
\end{table}

We see that in Russian, Ukrainian and Polish this noun is uninflectable, while inflecting in Czech and Serbo-Croat.\footnote{ \textrm{I treat Serbo-Croat (hbs) as a pluricentric language, like German or English, with four standards: Bosnian (bos), Croatian (hrv), Montenegrin (cnr), and Serbian (srp) following the 2017 ‘Deklaracija o zajedničkom jeziku’ (}\url{http://jezicinacionalizmi.com/deklaracija/}\textrm{). A linguistic overview can be found in \citet{CorbettBrowne2018}, and a good introduction to the complex issues of the sociolinguistic setting and language status is provided in Bugarski (2012, 2019).}} It is masculine in the languages in \tabref{tab:key:1}, except for Polish, where it is feminine. This noun points to a striking pattern: the languages form a ‘great semicircle’, running from Serbo-Croat in the southwest, where there are few uninflectable items, through Czech, Polish and Ukrainian, to Russian in the northeast, where there are many uninflectable items \citep[11-12]{Worth1966}. The picture is interestingly nuanced, as we shall see in \sectref{sec:key:3.2}.

The second family which will be well represented is Nakh-Dagestanian. These languages prove significant since on the one hand they show a remarkable degree of inflection, and on the other there are numerous instances of the inability to inflect. This Archi example is a good way in: 

Archi (aqc, \citealt{ChumakinaBond2016}: 112)\footnote{ \textrm{I follow the Leipzig Glossing Rules https://www.eva.mpg.de/lingua/resources/glossing-rules.php; [ ] indicates information that can be inferred from the use of the bare stem; < > marks infixation; ( ) is for inherent, non-overt feature values; bold is used as a flag to draw attention to relevant characteristics of examples. Abbreviations are the Leipzig ones with additions:} \textrm{\textsc{abs}}\textrm{: absolutive,} \textrm{\textsc{acc}}\textrm{: accusative,} \textrm{\textsc{dat}}\textrm{: dative,} \textrm{\textsc{decl}}\textrm{: declarative,} \textrm{\textsc{def}}\textrm{: definite,} \textrm{\textsc{dem}}\textrm{: demonstrative,} \textrm{\textsc{dimin}}\textrm{: diminutive,} \textrm{\textsc{erg}}\textrm{: ergative,} \textrm{\textsc{f}}\textrm{: feminine,} \textrm{\textsc{gen}}\textrm{: genitive,} \textrm{\textsc{ins}}\textrm{: instrumental,} \textrm{\textsc{loc}}\textrm{: locative,} \textrm{\textsc{m}}\textrm{: masculine,} \textrm{\textsc{med}}\textrm{: medial verb,} \textrm{\textsc{n}}\textrm{: neuter,} \textrm{\textsc{nmlz}}\textrm{: agent nominal,} \textrm{\textsc{nom}}\textrm{: nominative,} \textrm{\textsc{obj}}\textrm{: object,} \textrm{\textsc{pfv}}\textrm{: perfective,} \textrm{\textsc{pl}}\textrm{: plural,} \textrm{\textsc{prs}}\textrm{: present,} \textrm{\textsc{real}}\textrm{: realis,} \textrm{\textsc{refl}}\textrm{: reflexive,} \textrm{\textsc{sbj}}\textrm{: subject,} \textrm{\textsc{seq}}\textrm{: sequential,} \textrm{\textsc{sg}}\textrm{: singular. Features like} \textrm{\textbf{\textsc{number}}} \textrm{and values like} \textrm{\textsc{plural}} \textrm{are differentiated in this way where greater clarity is required.}} 

\ea%1
    \label{ex:key:1}
    \gll \\
         \\
    \glt
    \z % you might need an extra \z if this is the last of several subexamples

          d-ez  Ajša  \textbf{d-}akːu



  \textsc{sg.ii}{}-\textsc{1sg.dat}  Aisha(\textsc{ii)[sg.abs]}  \textbf{\textsc{sg.ii}}{}-see.\textsc{pfv}



  ‘I saw Aisha (woman’s name).’    [with agreement]


In \REF{ex:key:1} the verb inflects, agreeing with the absolutive argument \textit{Ajša} (a woman’s name). Nouns also inflect (here the stem \textit{Ajša} indicates the absolutive singular), and the pronoun \textit{dez}, which is first singular dative, also inflects to agree in number (\textsc{singular}) and gender (\textsc{ii}) with the absolutive argument. In fact, we find instances of inflection for all parts of speech. And yet, despite this prevalence of inflection, see example \REF{ex:key:2}: 

\ea%2
    \label{ex:key:2}
    \gll \\
         \\
    \glt
    \z % you might need an extra \z if this is the last of several subexamples

          d-ez  Ajša  kɬ’an



  \textsc{sg.ii}{}-\textsc{1sg.dat}  Aisha(\textsc{ii})[\textsc{sg.abs}]  love



  ‘I love Aisha.’  [no agreement]


In \REF{ex:key:2} we have the same construction as in \REF{ex:key:1}, namely a verb of emotion or perception, with the experiencer in the dative and the “object” in the absolutive. But in \REF{ex:key:2}, the verb does not agree. This is a lexical matter: some verbs agree (as in \REF{ex:key:1}) and some do not (as in \REF{ex:key:2}), and this extends to verbs of other semantic and syntactic types, which again include inflectable and uninflectable members. The pronoun still agrees in \REF{ex:key:2}, and more generally the syntax is ‘blind to the difference’ (\citealt{ChumakinaBond2016}: 112). This may disturb some of our assumptions, and so we shall return to Archi in \sectref{sec:key:3.1.2}.

\subsection{\textbf{Outline}}

I first define a canonical inflecting lexeme (\sectref{sec:key:2}). From this baseline I set up four scales of non\nobreakdash-canonicity. I take a canonical lexeme and compare outwards, measuring uninflectability \textit{across} items (\sectref{sec:key:3}). Then I compare inwards, measuring uninflectability \textit{within} items (\sectref{sec:key:4}). I analyse the context of uninflectability, looking at the conditions on it (the incoming conditions) in \sectref{sec:key:5}, followed by the external consequences (the outgoing impact) of uninflectability (\sectref{sec:key:6}). At a high level of generality these are the four theoretically possible dimensions of variability. I draw together the analyses in \sectref{sec:key:7}.

\section{\textbf{Canonical} \textbf{inflecting} \textbf{lexeme}}

I adopt a canonical approach (\sectref{sec:key:2.1}). I discuss relevant definitions (\sectref{sec:key:2.2}) and then consider the phenomena which border on uninflectedness (\sectref{sec:key:2.3}).

\subsection{\textbf{Canonical} \textbf{Typology}}

We wish to treat the phenomenon in a maximally inclusive way. Our method is well chosen:

‘the canonical approach breaks down complex concepts in a way that clarifies where disagreements may lie between different linguists and theoretical frameworks’ 

(\citealt{Nikolaeva2013}: 100; see also \citealt{Round2023} on this point)

In the canonical method, we analyse variable phenomena within and across languages to establish relevant scales (criteria). Then we push these scales to their logical end point, which gives us a baseline from which we can measure. This approach to measurement allows us ‘to conduct comparisons equally well within languages and between them’ (\citealt{RoundCorbett2020}). We identify more than one such scale, and this enables us to draw out the theoretical space to tease apart the clusters.

Canonical is a logical construct. It does not imply usual, normal, frequent, expected, unmarked or prototypical. As a comparison, zero kelvin is not a frequent or indeed desirable temperature, but it is useful as an agreed baseline from which we measure; similarly, canonical is a helpful starting point, a useful exploratory tool, not a claim about what we shall find.\footnote{For many examples of its usefulness see the Canonical Typology bibliography (permanent short-URL: tiny.cc/ctbib).} Among baselines, zero is favoured: 

‘... nothings are usually extremes. They tend to sit at one end of a spectrum. And when scientists want to explore a phenomenon they look for an extreme version of it, because the factors are often easier to spot.’   \citet[2]{Webb2013}

For uninflectability, the baseline from which we calibrate is canonical inflection. We start from a well-grounded analysis of the morphosyntactic features and their values for the given system (\citealt{Zaliznjak1973}, \citealt{Corbett2012}: 73-88). The feature values established ‘multiply out’, giving a paradigm with all the predicted paradigm cells. The lexical material is constant within lexemes and different across lexemes; conversely, inflectional material is different within lexemes and constant across lexemes. Thus every form of every lexeme will be distinct; and what can inflect will inflect \citep[150-151]{Corbett2015}. Hence we have a viable baseline, namely the canonically inflecting item, from which we can compare outwards across items.

\subsection{\textbf{Definitions} \textbf{(base)}}

What then is the phenomenon which we wish to investigate (the ‘base’ in Bond’s 2013 view)? We can state this simply: 

Uninflectability is the phenomenon where a member of an inflecting part of speech\footnote{I refer to ‘part of speech’, as is found particularly in HPSG. In LFG ‘lexical category’ is preferred, but ‘lexical’ brings unhelpful ambiguity, which is better avoided. And while some use ‘lexical category’ as a category of lexemes (equivalent to ‘part of speech’), others use it more narrowly for part of speech with lexical meaning, for example, noun, verb, and adjective; this is the prevalent usage in Minimalism.} fails to inflect. 

Compare Baerman, Brown \& \citet[27-35]{Corbett2005}, and \citet[142-143]{Spencer2020}. Note that the phrase ‘inflecting part of speech’ sets up the expectation, and we are interested in instances where that expectation is not met. These may be individual items, or they may form a subclass. As noted earlier, Spencer distinguishes ‘uninflectable’ items, which do not inflect despite an expectation that they will, from ‘uninflecting’ items, which belong to non\nobreakdash-inflecting classes, and hence there is no expectation that they will inflect \citep[142]{Spencer2020}. Note that, within an inflecting class, we are interested if most items or even all but one item inflect, and equally if most do not. The statistical question is useful, but it is orthogonal to other criteria. Given that the class of items is in principle inflecting, the fact that inflection is not found consistently means there is a phenomenon worthy of our attention. 

\subsection{\textbf{Closest} \textbf{related} \textbf{phenomena}}

Baerman, Brown \& \citet[27-35]{Corbett2005} are concerned with delimiting neutralization, syncretism and uninflectedness. They begin with two properties of canonical inflection. First, that feature values come together into fully distinguished combinations. A violation of this property is neutralization; here a distinction is not made in the inflectional morphology (for example, gender is not distinguished in the plural) and the syntax is also not sensitive to the distinction. The second property is that morphology always distinguishes syntactically relevant feature values. Violations of this property are uninflectedness, a total failure to distinguish relevant values, while syncretism (their main concern) represents a specific failure (for example, if dative and ablative case values are distinguished in the singular but not in the plural). To clarify the distinction, consider the data in \tabref{tab:key:2}.


\begin{tabularx}{\textwidth}{XXXXXXX}

\lsptoprule
%\hhline%%replace by cmidrule{~~-----}
\multicolumn{2}{c}{} & ‘boy’ & ‘book’ & ‘bone’ & ‘bog’ & ‘coat’\\
 \textsc{singular} & \textsc{nom} & mal′čik & knig-a & kost′ & bolot-o & pal′to\\
& \textsc{acc} & mal′čik-a & knig-u & kost′ & bolot-o & pal′to\\
& \textsc{gen} & mal′čik-a & knig-i & kost-i & bolot-a & pal′to\\
& \textsc{dat} & mal′čik-u & knig-e & kost-i & bolot-u & pal′to\\
& \textsc{loc} & mal′čik-e & knig-e & kost-i & bolot-e & pal′to\\
& \textsc{ins} & mal′čik-om & knig-oj & kost′-ju & bolot-om & pal′to  \\
 \textsc{plural} & \textsc{nom} & mal′čik-i & knig-i & kost-i & bolot-a & pal′to\\
& \textsc{acc} & mal′čik-ov & knig-i & kost-i & bolot-a & pal′to\\
& \textsc{gen} & mal′čik-ov & knig & kost-ej & bolot & pal′to\\
& \textsc{dat} & mal′čik-am & knig-am & kostj-am & bolot-am & pal′to\\
& \textsc{loc} & mal′čik-ax & knig-ax & kostj-ax & bolot-ax & pal′to\\
& \textsc{ins} & mal′čik-ami & knig-ami & kostj-ami & bolot-ami & pal′to  \\
\lspbottomrule
\end{tabularx}
%%please move \begin{table} just above \begin{tabular
\begin{table}
\caption{Illustrative nouns from Russian\footnote{ \textrm{These nouns illustrate the point we need, and are not fully representative of the system. For simplicity, stress alternations have not been included.} }}
\label{tab:key:2}
\end{table}

These nouns provide evidence to distinguish the core six case values of Russian. None of these nouns is fully canonical. We see several instances of syncretism. For example, \textit{kost′} ‘bone’ and \textit{boloto} ‘bog’ do not distinguish accusative from nominative in the singular or the plural, while \textit{kniga} ‘book’ has this syncretism only in the plural. \textit{Mal′čik} ‘boy’ has accusative-genitive syncretism in both numbers. Overall, \textit{kost′} ‘bone’ has fewer distinct forms than the other inflectable nouns, yet clearly it inflects. With smaller paradigms, the syncretism-uninflectability border becomes less clear; for discussion, with respect to Italian, see Thornton (2019a: 71-72, 90-91, and for the diachronic dimension see \citealt{Loporcaro2011}). Also in \tabref{tab:key:2} note the nominative singular of \textit{mal′čik} ‘boy’, and the genitive plural of \textit{kniga} ‘book’ and of \textit{boloto} ‘bog’. Here we find the zero form, the bare stem, what Stump (2017: 428n7) calls a ‘significative absence’. There is no affixal inflection but, within the system, the morphosyntactic specification of forms with no overt inflection is as clearly identifiable as for the other forms; such forms do not count as uninflectable. In contrast, the noun \textit{pal′to} ‘coat’ is uninflectable. This noun has a full paradigm, in the sense that if can appear in all syntactic slots, but the different morphosyntactic requirements are all met by the same form. This uninflectable noun is not distinct syntactically from inflectable nouns; for instance, it controls number agreement just as normal inflecting nouns do. This is shown in these contrasting examples \REF{ex:key:3}, appropriate for nominative or accusative cases, which occur several times in the Russian National Corpus (\href{http://www.ruscorpora.ru}{{www.ruscorpora.ru}}{)}:

Russian (rus)

\ea%3
    \label{ex:key:3}
    \gll \\
         \\
    \glt
    \z % you might need an extra \z if this is the last of several subexamples

        a  zimn-ee  pal′to  \REF{ex:key:3}b  zimn-ie  pal′to

  winter-\textsc{sg.n}  coat    winter-\textsc{pl}  coats

  ‘(a) winter coat’  ‘winter coats’

Baerman, Brown \& Corbett define uninflectedness as follows (2005: 33):

\begin{itemize}
\item \begin{styleListParagraph}
There is, in certain lexemes only, a loss of all values of a particular feature F found elsewhere in the language. This loss may depend on the presence of a particular combination of values of one or more other features (the context).
\end{styleListParagraph}
\item \begin{styleListParagraph}
Other syntactic objects distinguish values of feature F, either generally or in the given context, and feature F is therefore syntactically relevant.
\end{styleListParagraph}
\end{itemize}

The key point is that uninflectedness, and more specifically uninflectability, is about morphology being unresponsive to a feature that is syntactically relevant, and this definition allows for partial uninflectability.

The other ‘boundary’ to consider is that with defectiveness. In principle, the distinction between uninflectability and defectiveness is clear. Defectiveness means that the required form is missing and hence that an anticipated morphosyntactic specification cannot be realized (\citealt{BaermanCorbett2010}, Chuang, Brown \citealt{BaayenEvans2022}, \citealt{Sims2023}). Uninflectability is less dramatic: the morphosyntactic specification is not overtly realized, but the uninflected form is available and there is no resulting syntactic problem. Table~3 shows different possibilities. 


\begin{tabularx}{\textwidth}{XXXXXX}

\lsptoprule
%\hhline%%replace by cmidrule{~~----}
\multicolumn{2}{c}{} & ‘dream’ & ‘coat’ & ‘riding breeches’ & ‘scissors’\\
 \textsc{singular} & \textsc{nom} & mečt-a & pal′to & *** & ***\\
& \textsc{acc} & mečt-u & pal′to & *** & ***\\
& \textsc{gen} & mečt-y & pal′to & *** & ***\\
& \textsc{dat} & mečt-e & pal′to & *** & ***\\
& \textsc{loc} & mečt-e & pal′to & *** & ***\\
& \textsc{ins} & mečt-oj & pal′to & *** & ***\\
 \textsc{plural} & \textsc{nom} & mečt-y & pal′to & galife & nožnic-y\\
& \textsc{acc} & mečt-y & pal′to & galife & nožnic-y\\
& \textsc{gen} & *** & pal′to & galife & nožnic\\
& \textsc{dat} & mečt-am & pal′to & galife & nožnic-am\\
& \textsc{loc} & mečt-ax & pal′to & galife & nožnic-ami\\
& \textsc{ins} & mečt-ami & pal′to & galife & nožnic-ax\\
\lspbottomrule
\end{tabularx}
%%please move \begin{table} just above \begin{tabular
\begin{table}
\caption{Defectiveness and uninflectability illustrated from Russian}
\label{tab:key:3}
\end{table}

For many Russian speakers, it is not possible to form the genitive singular of the noun \textit{mečta} ‘dream’. There is a gap here: the noun is defective. This contrasts with uninflectable \textit{pal′to} ‘coat’, which has an unproblematic genitive plural, one which is like all its other forms. However, the contrast between defectiveness and uninflectability is not always so obvious. A noun may be uninflectable and also defective. We see this with \textit{galife} ‘riding breeches’ (Table~3), which does not inflect, and controls only plural agreement:

Russian \citep[60]{Corbett2019}

\ea%4
    \label{ex:key:4}
    \gll \\
         \\
    \glt
    \z % you might need an extra \z if this is the last of several subexamples

          v  ser-yx  galife

  in  grey\textsc{{}-pl.loc}  riding.breeches

  ‘in grey riding breeches’  (Ju. Trifonov, \textit{Dom na naberežnoj} 1976)

This noun\footnote{ \textrm{See \citealt{Dunn1989} for further uninflectable pluralia tantum nouns.}} is to be compared with nouns like \textit{nožnicy} ‘scissors’ (\tabref{tab:key:3}), which is also a plurale tantum noun (hence defective in having no singular), but it inflects fully in the plural \citep[59-60]{Corbett2019}. Given that defectiveness and uninflectability may intersect, we might still hope to maintain a clear distinction: uninflectability brings no syntactic problem while defectiveness causes the syntax to crash. But even here there are complications. Items may vary, being uninflectable or defective, according to other feature values, as we shall see in \sectref{sec:key:6.2.1} and \sectref{sec:key:6.2.2} (and compare \citet{NikolaevBermel2022}).

\subsection{\textbf{A} \textbf{specific} \textbf{issue} \textbf{with} \textbf{uninflectable} \textbf{nouns:} \textbf{gender} \textbf{assignment}}

Recall from \tabref{tab:key:1} that \textit{kakadu} ‘cockatoo’ in a selection of Slavonic languages varies according to whether it inflects, and according to its gender. There is a large issue here, but one restricted to nouns, namely gender assignment. It is important not to simplify: there is often a link between borrowings and uninflectability (but not all borrowings are uninflectable, and not all uninflectable nouns are borrowings); furthermore, not all uninflectable nouns need any additional specification as to gender. Indeed, the null hypothesis is that uninflectable nouns behave as ordinary nouns for the purposes of gender assignment. That is, assignment will be via semantic similarity, or formal similarity (whether morphological or phonological), see Fedden, Guzmán Naranjo \& Corbett (to appear). However, the problem that prevents a noun from inflecting may equally be an issue for gender assignment. Thus, uninflectability here is important for what it can tell us about gender assignment, and in the case of borrowings, about the degree to which borrowed items are integrated. There is a considerable literature. The problem is acute in languages like those of the Slavonic family, where inflectional class is a strong predictor of gender (except when semantic assignment of animate nouns takes precedence). Within Slavonic, Russian is well studied, as in \citet{Murphy2000}, \citet{Wang2014}, \citet{Matushansky2022}, Chuprinko, \citet{MagomedovaSlioussar2023} and \citet[148-158]{Timberlake2004}, who provides a fine summary; for comparisons across the family, see \citet[193-197]{Thomas1983} and Doleschal (2000-2001: 48-151, 2025). 

\section{\textbf{Extent} \textbf{of} \textbf{uninflectability} \textbf{across} \textbf{items}}

To make progress we need to measure, and we start by measuring across lexemes (\sectref{sec:key:3.1}). Then we ask how we can specify which items inflect and which are uninflectable (\sectref{sec:key:3.2}).

\subsection{\textbf{Measuring} \textbf{across} \textbf{items}}

How many uninflectable items are there? In the canonical case, clearly none. We are interested in the extent of the difference between the canonical zero and what we find, and we have three types of comparison. The key type is within the part of speech, which sets up the expectation of inflectability (\sectref{sec:key:3.1.1}). We can also compare more generally within the language (\sectref{sec:key:3.1.2}) and even across languages (\sectref{sec:key:3.1.3}). We can then approach the question of which items fail to inflect (\sectref{sec:key:3.2}). 

\subsubsection{\textbf{Comparison} \textbf{within} \textbf{the} \textbf{part} \textbf{of} \textbf{speech} \textbf{(the} \textbf{key} \textbf{comparison)}}

We start from the baseline expectation that behaviours will line up; specifically, members of a part of speech, defined in syntactic terms, will have the same morphological behaviour (\citealt{Spencer2005}: 101, \citealt{Corbett2013}). This is part of the general approach in Canonical Typology, namely that our criteria converge (\citealt{BrownChumakina2013}: 9, \citealt{CorbettFedden2016}: 527). This would mean that for our current purposes, within a given part of speech, all members inflect, or none do. We are interested in the instances where that expectation is not met. And, importantly, we are equally concerned with those examples where uninflectability is the minority behaviour and those where inflecting is unusual. We may be more familiar with the Russian-type situation, that is, where most nouns inflect, but some are uninflectable (like \textit{pal′to} ‘coat’ in \tabref{tab:key:3}); however, the situation where a minority of items inflects, as in \tabref{tab:key:4} below, is just as interesting. One reason for this is that – in most current theories of syntax – mechanisms for distributing and constraining feature values (for government and agreement, however these are modelled) apply ‘blindly’; they do not ‘peek ahead’ to check whether the particular governee or agreement target has distinct forms for the feature values distributed. In other words, syntax is ‘morphology-free’ (\citealt{Zwicky1996}: 301, discussed in \citealt{Corbett2009}). Where only a (small) minority of the potential items is involved, this is relevant both for the functioning of the syntactic system and for its acquisition (see, for instance, \citealt{ChumakinaBond2016}: 111-117). 

\subsubsection{\textbf{Comparison} \textbf{within} \textbf{the} \textbf{language} \textbf{(secondary)} }

There are languages where comparison across the parts of speech yields interesting results. For example, Serbo-Croat nominals typically retain inflection, yet numerals have largely lost it. This forms a surprising contrast to Russian, which has many uninflected nouns, yet retains inflection for numerals. 

In Archi (introduced in \sectref{sec:key:1.3}) we find another unexpected contrast. On the one hand, we find inflection in all parts of speech. Example \REF{ex:key:5} illustrates part of the picture:

Archi (aqc, \citealt{Kibrik1994}: 349): inflection in different parts of speech

\ea%5
    \label{ex:key:5}
    \gll \\
         \\
    \glt
    \z % you might need an extra \z if this is the last of several subexamples

          buwa-mu  \textbf{b-ez}  ditːa<b>u  χːʷalli  a<b>u

mother(\textsc{ii})-\textsc{sg.erg}  \textbf{\textsc{sg.iii}}{}-1\textsc{sg}.\textsc{dat}  early<\textsc{sg.iii}>  bread(\textsc{iii)[sg.abs]}  make<\textsc{sg.iii}>\textsc{pfv}

  ‘Mother made bread for me early.’

In \REF{ex:key:5} we see case and number inflection on the noun \textit{buwa} ‘mother’. This example also shows other parts of speech inflecting for agreement (with the absolutive argument \textit{χːʷalli} ‘bread’). Indeed there is agreement inflection found on all parts of speech apart from nouns (\citealt{ChumakinaCorbett2015}: 94-95); in addition, for some parts of speech, there is inflection for other features too.\footnote{\textrm{It is worth noting that Archi has instances of ‘frivolous infixation’, not determined by phonology (\citealt{ChumakinaCorbett2015}).}} The contrast is that within these inflecting parts of speech there may be many items which do not inflect fully. This is clear from \tabref{tab:key:4}.


\begin{tabularx}{\textwidth}{XXXX}

\lsptoprule
%\hhline%%replace by cmidrule{~---}
\multicolumn{1}{c}{~} & \raggedleft total & \raggedleft agreeing & \raggedleft \% agreeing\\
verbs & \raggedleft 1248 & \raggedleft 399 & \raggedleft 32.0\\
adverbs & \raggedleft 383 & \raggedleft 13 & \raggedleft 3.4\\
postpositions & \raggedleft 34 & \raggedleft 1 & \raggedleft 2.9\\
discourse clitics/particles & \raggedleft 4 & \raggedleft 1 & \raggedleft 25.0\\
\lspbottomrule
\end{tabularx}
%%please move \begin{table} just above \begin{tabular
\begin{table}
\caption{Lexemes which inflect for agreement in the Archi dictionary (\citealt{ChumakinaCorbett2015}: 95)\footnote{ \textrm{Pronouns are not included in \tabref{tab:key:4}, since we cannot simply count lexemes in this instance, as we shall see in \sectref{sec:key:4.1}. \citet{Chumakina2023} reports on adverbs which inflect for orientation case markers (such as allative ‘towards’); there are about ten such adverbs in Archi, but no adverb inflects for orientation case as well as for gender and number (Marina Chumakina, personal communication 13.11.2023).}}}
\label{tab:key:4}
\end{table}

At this stage, the key point about Archi is the existence of parts of speech where full inflection is found with only a minority of items (and compare \citealt{Palancar2025} on Amuzgan). We shall see another example in Mian in \sectref{sec:key:4.2}, particularly \tabref{tab:key:7}. It might be suggested that minority items which inflect (as in \tabref{tab:key:4}) could be treated as instances of overdifferentiation, but they are more ‘serious’ than that. Overdifferentiation involves additional feature values, as when Kolami numerals show three gender values as opposed to the two values found with other parts of speech (\citealt{Emeneau1955}: 56, \citealt{Corbett1991}: 168). In other words, the feature gender is common to numerals and other parts of speech, but numerals have an additional gender value. In \tabref{tab:key:4}, it is a matter of inflecting for a feature as a whole, or not, which is a more radical difference than indicated by the term overdifferentiation. We return to Archi in \sectref{sec:key:4.2}, to take further the point about ‘full’ inflection, since items can fail to show gender and number agreement, and still inflect for other features. 

\subsubsection{\textbf{Comparison} \textbf{cross-linguistically} }

We have general cross-linguistic (typological) expectations: we are not surprised to find verbs which inflect, while we may be relatively confident that conjunctions do not. And yet even this surprising possibility can be found:

Walman (van, Torricelli, New Guinea, \citealt{BrownDryer2008}: 534)

\ea%6
    \label{ex:key:6}
    \gll \\
         \\
    \glt
    \z % you might need an extra \z if this is the last of several subexamples

          w-ru  chuto  rounu  alpa  \textbf{w-aro-l}  nyanam

  \textsc{gen-3sg.f}  female  old  one[\textsc{f}]  \textbf{\textsc{3sg.f}}\textsc{{}-}and\textsc{{}-}\textbf{\textsc{3sg.dimin}}  child

  nngkal  ngo-l  pa  y-an  nakol

  small  one-\textsc{dimin}  \textsc{dem}  \textsc{3pl}{}-be.at  village

  ‘(Only) one old woman and a small child were left behind in the village.’

Observe how \textit{{}-aro-} ‘and’ agrees with both conjuncts: it has the third singular prefix \textit{w-} to agree with the first conjunct and the diminutive suffix \textit{{}-l} to agree with the second (just as Walman verbs can mark subject and object). It is morphologically a verb (and can be used as a verb with the sense ‘be with’), though syntactically it can behave like a conjunction. This shows that the typological picture is not clear; we know in principle that inflection is more frequent in certain parts of speech than in others, but we lack a full typological survey.

\subsection{\textbf{Determining} \textbf{which} \textbf{items} \textbf{(fail} \textbf{to)} \textbf{inflect}}

Naturally we ask why some items fail to inflect (or indeed why some items inflect). We may frame the question in terms of ‘two ancient dichotomies’: we ask what is unpredictable as opposed to predictable/rule-governed, and what is arbitrary as opposed to functionally motivated/natural \citep[107]{Bye2015}. For the first dichotomy, unpredictable as opposed to rule governed, there are several factors adduced as determining whether or not a particular item will inflect (given that it belongs to a part of speech where inflection is found); \citet{Fedden2019} points to these: phonotactic, phonological, morphological and semantic (and compare \citealt{Doleschal2000}-2001: 17-19 for systematic vs unsystematic uninflectability). With the part of speech noun in Slavonic languages, Doleschal (2000-2001: 48-151) considers the specifics of personal names, geographical names,\footnote{ \textrm{Russian personal names are well described in \citet[153-158]{Timberlake2004}; see also the discussion in Spencer (this volume). The suggestion that in World War II, Russian official sources did not inflect place names to avoid confusion between similar names goes back to \citet[105]{Mirtov1953}.} } the names of organizations and products, foreign words and abbreviations. And specifically in a study of Archi (see \tabref{tab:key:4} above), \citet{ChumakinaCorbett2015} show the need for reference to phonology, derivation, and morphological type (e.g. dynamic vs stative verbs) to determine which items inflect. 

The second dichotomy contrasts arbitrary and functionally motivated (compare \citealt{Spencer2020}: 147). A recurring and convincing type of example here involves items where the expected inflection would bring the item into the wrong class. Take borrowed female names like \textit{Dolôres} in Serbo-Croat \citep[378]{Browne1993a}. This name has the form of a bare stem and so “ought” to inflect in the same way as nouns like \textit{Jòvan}, a man’s name. But this would make it like a masculine noun, and so it remains uninflectable.\footnote{ \textrm{There is another theoretical alternative, the third inflection class, where nouns like} \textrm{\textit{svȃst} }\textrm{‘sister-in-law’ would be a potential model. This class includes few animates and does not attract borrowed names.} } A related type is that where the motivation is better alignment of the item with a part of speech or subtype within it. In German, for instance, proper names of different types have increasingly formed a subclass of nouns, reducing their degree of inflection when compared with the part of speech as a whole (see \citealt{Nübling2017}: 344-349). For Italian, D’\citet{AchilleThornton2003} and Thornton \& D’\citet{Achille2025} track the rise of uninflectability over several centuries, including some nouns which do not fit the main correspondences between inflection and gender. For Polish see \citet{Gaszewski2025} and for Japanese see \citet{Köhlich2025}.

It is here that Slavonic languages are particularly apposite. Recall from \sectref{sec:key:1.3} Worth’s (1966: 11-12) point that the languages form a ‘great semicircle’, running from Serbo-Croat, where there are few uninflectable items, through Czech, Polish and Ukrainian, to Russian, where there are more uninflectable items.\footnote{\citet{Naylor1982} suggests a phonological explanation for why Serbo-Croat has so few uninflectable nouns, namely the presence of vowel clusters, which offers a route for assimilating loans with a shape that would be problematic for Russian. However, many of the uninflectable nouns of Russian have a phonological shape which could fit readily into the inflectional system.}  The situation is more nuanced in that, for instance, Russian has many uninflectable nouns, yet retains inflection for numerals, while Serbo-Croat is strict on inflecting nouns but has developed uninflectable numerals (as we noted above, and to be discussed in \sectref{sec:key:6.2.2}). We might expect that borrowings would fail to inflect, if their phonological shape did not fit with the inflectional system, and that they might be accommodated over time; hence that the underlying trend would be for uninflectables to become inflectable. But there is clear evidence of instances of borrowings initially being inflected, and only later becoming uninflectable. In Russian in particular, nouns which were borrowed (in many instances from French) were in some instances inflected, but later were established as uninflectable. \citet[161-164]{Mučnik1964} shows how in 18\textsuperscript{th} century Russian, borrowed nouns tended to be inflected, and that this tendency was reversed during the 19\textsuperscript{th} century, with the uninflectable status of many nouns becoming firmer in the early 20\textsuperscript{th} century. This is the result of the pressure of educated speakers of the standard language. A convenient summary with sources is given by Comrie, Stone \& \citet[117-123]{Polinsky1996}; see \citet[45-55]{Panov1968} for speakers’ negative reactions when asked about inflecting such nouns; there is useful discussion in \citet{Dunn1988} and a wealth of data in Doleschal (2000-2001).\footnote{It is not only Russian which shows an increase in uninflectable items. \citet{Auty1968} gives evidence to suggest a partly similar development in Czech.} This discussion shows that, more generally, to demonstrate what the determining factors are, diachronic evidence can prove valuable; for helpful use of diachronic data see further Copot, Da Silva, Beji, \citet{WatiezBonami2025}, \citet{Friedewald2025}, \citet{RennerRenwick2025} and \citet{Souag2025}.

\section{\textbf{Extent} \textbf{of} \textbf{uninflectability} \textbf{within} \textbf{items}}

We turn to uninflectability within items. Here we might anticipate that uninflectability would be a black and white affair: lexemes inflect or they do not. But the situation is more subtle. Paradigms may be split into an inflectable and an uninflectable portion (or segment); in such an instance we need to define the split (\sectref{sec:key:4.1}).\footnote{ \textrm{There are instances where subclasses of lexemes lack full paradigms. Examples included pluralia tantum nouns, location nominals and so on (Baerman, \citealt{BrownCorbett2017}: 71-74, 79). Such items have a paradigm structure different from that of the main class (they lack certain cells) and are not our concern here. We are concerned with instances where, given the paradigm, some cells are filled but lack inflection.} } There are also situations where the system of feature specification may need elaboration: we need to check whether a simple specification is adequate or a multiple feature specification is justified (\sectref{sec:key:4.2}). These two criteria allow us to build a full typology (\sectref{sec:key:4.3}). 

\subsection{\textbf{Definition} \textbf{of} \textbf{inflectability} \textbf{splits}}

Consider the Polish (pol) data in \tabref{tab:key:5}. 


\begin{tabularx}{\textwidth}{XXXXX}

\lsptoprule
%\hhline%%replace by cmidrule{~~---}
\multicolumn{1}{c}{} &  & \textit{wino} ‘wine’ & \textit{muzeum} ‘museum’ & \textit{hobby} ‘hobby’\\
 \textsc{singular} & \textsc{nominative} & wino & muzeum & hobby\\
& \textsc{accusative} & wino & muzeum & hobby\\
& \textsc{genitive} & wina & muzeum & hobby\\
& \textsc{dative} & winu & muzeum & hobby\\
& \textsc{locative} & winie & muzeum & hobby\\
& \textsc{instrumental} & winem & muzeum & hobby\\
 \textsc{uuuplural} & \textsc{nominative} & wina & muzea & hobby\\
& \textsc{accusative} & wina & muzea & hobby\\
& \textsc{genitive} & win & muzeów\footnotemark{} & hobby\\
& \textsc{dative} & winom & muzeom & hobby\\
& \textsc{locative} & winach & muzeach & hobby\\
& \textsc{instrumental} & winami & muzeami & hobby\\
\lspbottomrule
\end{tabularx}
\footnotetext{ The genitive plural ending is not what we would expect (namely the bare stem, as with \textit{win}, the genitive plural of \textit{wino} ‘wine’). The \textit{{}-ów} inflection is found with the main inflection class of masculine nouns and with some pluralia tantum nouns. Note also that according to \citet[120]{Swan2002} Polish uninflectable nouns do sometimes inflect.}
%%please move \begin{table} just above \begin{tabular
\begin{table}
\caption{Three Polish nouns (\citealt{Swan2002}: 116, 118, 121)}
\label{tab:key:5}
\end{table}

Polish has straightforward, inflecting nouns, like \textit{wino} ‘wine’ in \tabref{tab:key:5}, and there are different patterns (inflection classes); it also has uninflectable nouns like \textit{hobby} ‘hobby’. But Table~5 shows an instance which falls between, namely \textit{muzeum} ‘museum’. The paradigm of this noun is split by uninflectability. In the plural we have full inflection (for number and case), while in the singular there is a different stem, and uninflectability. We could in principle measure from either extreme: that is, the noun inflects (apart from the singular) or it is uninflectable (apart from the plural). For consistency, we have been measuring from full inflection. \textit{Muzeum} ‘museum’ is an unusual instance of heteroclisis, in that nouns like this take their forms from two different inflection classes (one being uninflectable), as indicated by the colour coding. We see that the uninflectable portion consists of a slab of the paradigm, that is, in a motivated portion (hence we have featural uninflectability). 

The majority of Polish nouns inflect (whether like \textit{wino} ‘wine’ or according to some other pattern); a minority are uninflectable (like \textit{hobby} ‘hobby’), and a few behave like \textit{muzeum} ‘museum’, as listed by \citet[118]{Swan2002}. However, for the purposes of the current section that is interesting but secondary information. Here we abstract away from the number of items which behave in a particular way (see \sectref{sec:key:3.1.1}). The key point is that there is more going on than a simple difference between fully inflecting and not inflecting. Nouns like \textit{muzeum} ‘museum’ suggest that a single dimension (inflecting – uninflectable) will not be adequate. As Brown \& Chumakina point out more generally (2013: 6):

... the space defined by the criteria may contain dimensions which differ in their nature (i.e. in terms of the number of points in the dimension, and whether the dimension is discrete or gradient).

We therefore need to investigate further. We noted that the uninflectable portion of nouns like \textit{muzeum} ‘museum’ consists of a slab of the paradigm, in this instance all the singular values. We ask whether partial uninflectability has to be this way, or whether the logical possibility of uninflectedness according to particular cells, can be found. This distinction is valid elsewhere. Thus pluralia tantum nouns are defective for a slab of the paradigm; such motivated portions are sometimes called ‘sub-paradigms’. And this is a dimension that can contribute to our typology (as we see later in \figref{fig:key:1}).

There is a suggested instance, the Russian noun \textit{èxo} ‘echo’. \citet[133-134]{Unbegaun1947} states that this noun is more likely to inflect in the instrumental singular (\textit{èxom}) than elsewhere. \citet[182-183]{Thomas1983} goes further, saying that this noun inflects \textit{only} in the instrumental singular (and gives references).\footnote{The original claim was made many years ago, and of course change is likely (as with other items cited from earlier periods); see Comrie, Stone \& \citet[120]{Polinsky1996} for sources of the 1960s which suggested that \textit{èxo} ‘echo’ was less frequently inflected than it had been previously. Our concern is with what is possible, even if a previous situation no longer holds. See \citet{Esher2025} for a complex series of incremental changes in French, correlated with inflectional class, phonological shape and gender; the loss of case contrasts involves regular sound change, analogy and syntactic change, while the loss of number contrasts is principally due to sound change.} We should reflect on this claim: a given lexeme can inflect or not inflect according to the specific cell in the paradigm. If correct, this claim implies a dramatic expansion of the notion ‘possible lexeme’.\footnote{ \textrm{Items can be defective cell by cell, as illustrated for Latin by \citet[229]{Rhodes1987}; thus} \textrm{\textit{vīs}} \textrm{‘force’ lacks the genitive and dative singular. The suggestion we are discussing is “worse” in the sense that a given item is claimed to be non-defective, yet not to inflect in all cells of the paradigm.}} 

To test this, I searched the Russian National Corpus (available at \href{http://www.ruscorpora.ru}{{www.ruscorpora.ru}}), taking the automatically disambiguated texts only, which yielded over 5000 potential examples. There were extremely few plurals, and just the singular examples were retained.\footnote{\textrm{Examples from two authors were omitted: one had the noun in a different inflection class, with nominative} \textrm{\textit{èxa}}\textrm{, the other had} \textrm{\textit{Èxe}} \textrm{as a character’s name.}} The parser allows us to establish, for each cell, how many times the noun is inflected, and how often it remains uninflected. Compare the following examples, both involving the instrumental:

\ea%7
    \label{ex:key:7}
    \gll \\
         \\
    \glt
    \z % you might need an extra \z if this is the last of several subexamples

          ... ee  golos  razmeta-l-sja  po  perexod-u

      her  voice(\textsc{m)[sg.nom}]  disperse-\textsc{pst[sg.m}]\textsc{{}-refl}  round  passage-\textsc{sg.dat}

  slab-ym  \textbf{èxo}

  weak-\textsc{sg.ins.n}  echo

  ‘her voice scattered around the passage as a weak echo’ 

    (Evgenij Proškin, \textit{Mexanika večnosti} 2001)

In \REF{ex:key:7} the syntax and the agreeing adjective make clear that the case is instrumental, and the noun \textit{èxo} ‘echo’ remains uninflected. Now compare \REF{ex:key:8}:

\ea%8
    \label{ex:key:8}
    \gll \\
         \\
    \glt
    \z % you might need an extra \z if this is the last of several subexamples

          Moj  golos  gulk-im  \textbf{èx-om}  

  my[\textsc{sg.nom.m}]  voice(\textsc{m)[sg.nom}]  \textstyleyiiiqfc{resonant}{}-\textsc{sg.ins.n}  echo-\textsc{sg.ins}

  raznës-sja  po  zal-e ...

  spread.\textsc{pst[sg.m}]\textsc{{}-refl}  round  hall-\textsc{sg.dat}

  ‘\textstyleyiiiqfc{My voice spread as a resonant echo through the hall ...}’ 

    (Evgenija Pankratova, \textit{Angel čerdačnogo okna} 2015)

Example \REF{ex:key:8} is similar to \REF{ex:key:7}, but with an inflected form of \textit{èxo} ‘echo’. As we shall see, the proportions do indeed vary from cell to cell, with the instrumental showing the highest proportion of inflected forms. Here we hit a problem. The number of instances of the use of of \textit{èxo} ‘echo’ in the instrumental case is suspiciously high, but we need a standard of comparison. We need to know what the distribution of a typical Russian noun looks like. Progress towards this goal is recorded in \citet{Koptev2008}; he gives data on case but does not distinguish according to number. As a step towards this, I offer for comparison the frequency ranking of the singular case values of nouns in three corpora: first a corpus of 20\textsuperscript{th} century Russian fiction \citep{Sharoff2006}. Second, a web corpus (${\sim}$100 million words) from 2013 (many thanks to Dunstan Brown for his help here). Third, counts from Slioussar \& Samoilova (2015: \tabref{tab:key:6}), also using the Russian National Corpus (the grammatically disambiguated subcorpus SynTagRus yielding {\textasciitilde}1.2 million singular forms). These three corpora allow us to derive comparative data on individual case values, to compare with those of \textit{èxo} ‘echo’, as in Table~6.


\begin{tabularx}{\textwidth}{XXXXXXX}

\lsptoprule

Russian \textit{èxo} ‘echo’ \textsc{singular} & examples (RNC) & \% inflected & rank by total & \multicolumn{3}{c}{ Rank for typical noun}\\
&  &  &  & Corpus 1 & Corpus 2 & Corpus 3\\
\textsc{nominative} & \raggedleft 2996 & +/- & 1 & 1 & 2 & 1\\
\textsc{accusative} & \raggedleft 444 & +/- & 3 & 2 & 3 & 3\\
\textsc{genitive} & \raggedleft 274 & 86.9 & 4 & 3 & 1 & 2\\
\textsc{dative} & \raggedleft 98 & 80.6 & 5 & 6 & 6 & 6\\
\textsc{locative} & \raggedleft 63 & 66.7 & 6 & 4 & 4 & 4\\
\textsc{instrumental} & \raggedleft 1138 & 98.5 & 2 & 5 & 5 & 5\\
\lspbottomrule
\end{tabularx}
%%please move \begin{table} just above \begin{tabular
\begin{table}
\caption{Optional inflection: Russian \textit{èxo} ‘echo’ (rank of singular only) vs typical noun}
\label{tab:key:6}
\end{table}

Notes:  RNC: Russian National Corpus 

Corpus 1: 20\textsuperscript{th} century fiction: rank order from Brown, Tiberius \& \citet[522]{Corbett2007} 

Corpus 2: web corpus: Araneum Russicum III Minus

Corpus 3: RNC, Slioussar \& Samoilova (2015: \tabref{tab:key:6}) 

\tabref{tab:key:6} first gives the number of examples of \textit{èxo} ‘echo’ found in the Russian National Corpus and then the portion of those which are inflected. For instances which appear in syntactic positions requiring nominative or accusative, we cannot tell whether the noun is inflected or not, given that the form \textit{èxo} ‘echo’ is a fine nominative or accusative singular, hence the indicator +/-. For the remaining case values, we see that the noun is more often inflected than not. This is not as expected, but it is still directly relevant: this noun is neither simply inflecting nor uninflectable but rather it is partially inflecting. And indeed the instrumental is the case value where we find the highest degree of inflection (98.5\%). There is an interesting pattern: as the instances of a particular case value increase, so does the likelihood of inflection. And this leads to the question of what the expected distribution of frequencies is. For this we look to other corpora where a partial answer can be obtained. For each I give the rank order of the singular case values (\tabref{tab:key:6}). In all three corpora we find the instrumental as less frequent than all but the locative (for singular nouns as a whole). There is indeed something unusual about the noun \textit{èxo} ‘echo’: the frequency distribution of its case values does not fit the general pattern of nouns,\footnote{ \textrm{There is clearly more to be done to understand the typical pattern. For instance, Slioussar \& Samoilova (2015: \tabref{tab:key:6}) compare animate and inanimate nouns. Here the relative frequencies of the different cases vary considerably; however, the instrumental singular for inanimate nouns like} \textrm{\textit{èxo}} \textrm{‘echo’ again ranks fifth.} } since the instrumental is its most frequent case value, apart from the nominative. It does inflect, but not fully (\citealt{Unbegaun1947} was partially right), and we are interested in partial inflection, whatever the distribution of inflection vs lack of inflection. Moreover, the likelihood of inflection for \textit{èxo} ‘echo’ depends on the case value involved – that is the key point for our typology. 

I draw three points from the discussion of Russian \textit{èxo} ‘echo’. First, since it (frequently) inflects, but may also not inflect, this means that we find instances of overabundance \citep{Thornton2019b}: we have two forms available to realize the same cell in the paradigm. Other instances of this kind have been noted, in Czech \citep[12]{Auty1968}, Polish (\citealt{Swan2002}: 120, Krzyżanowski 2020: 53), and Belarusian (\citealt{Mayo1993}: 905 on \textit{maci} ‘mother’).

Second, Russian \textit{èxo} ‘echo’ shows how idiosyncratic a noun can be in inflectional terms, specifically in terms of the impact of case values on inflectability. The interaction of different types of inflectional irregularity was highlighted in \citet{Corbett2011}. Further, we know that there are nouns which have special distributions with respect to number. We find number dominance (for instance, English \textit{lips} occurs much more frequently in the plural than in the singular, while the usual distribution is for the singular to be more common). There is also number orientation, whereby nouns have singular and plural forms available, but do not use them in an unconstrained way (for instance, Russian \textit{ogurcy} ‘cucumbers’ has a regular singular, but there are contextual limitations on the use of this singular). For examples and sources on number dominance and number orientation see \citet[67-71]{Corbett2019}. And there can be very specific phonetic properties of given items (see \citealt{SchlechtwegCorbett2023}: 69-70 for references). All this shows that there is much more to be understood about the profile of individual lexical items, including their propensity to inflect.

Third, and most important for current purposes, inflectability can vary according to individual paradigm cells. We have seen this for Russian \textit{èxo} ‘echo’; \citet{Doleschal2002a} suggests that the case value has an effect on inflectability in Czech, and further evidence of optional inflection, from Bulgarian and Serbian dialects, is provided by \citet{KyusevaEtAl2025}. And more generally, there is more to paradigm cells than their featural specification (see Bonami, \citealt{KyjánekWauquier2023}: 317).

For current purposes, we see that by contrasting examples from Polish and Russian we have demonstrated an opposition according to the inflectable portion of the paradigm (see Figure~1).


\begin{tabularx}{\textwidth}{XX}

\lsptoprule

\multicolumn{2}{c}{ definition of inflectability split}\\
slab (featural) & cell(s) (morphomic)\\
Polish \textit{muzeum} ‘museum’ & Russian \textit{èxo} ‘echo’\\
\lspbottomrule
\end{tabularx}
\begin{figure}
\caption{Partial inflectability: examples of the first dimension of comparison}
\label{fig:key:1}
\end{figure}

As \figref{fig:key:1} indicates, lexemes may be split for inflectability into portions, which may be defined as a slab of the paradigm or by particular cells. 

\subsection{\textbf{System} \textbf{of} \textbf{feature} \textbf{specification}}

I now turn to the second, cross-cutting criterion; this concerns the feature system. Consider this example from Mian (Mountain Ok family, within Trans New Guinea):

Mian (mpt, \citealt{Fedden2022}: 284)

\ea%9
    \label{ex:key:9}
    \gll \\
         \\
    \glt
    \z % you might need an extra \z if this is the last of several subexamples

          máam=e  \textbf{a}{}-nâ’-n-\textbf{o}=be  

  mosquito=\textsc{sg.m}  \textbf{\textsc{3sg.m.obj}}\textbf{{}-}hit-\textsc{real-}\textbf{\textsc{3sg.f.sbj}}\textsc{=decl}

  ‘She hit the mosquito.’

The verb has two featural specifications for person, number and gender. We have a verb that marks agreement with the object by means of a prefix, agreement is in person (\textsc{3}), number (\textsc{sg}) and gender (\textsc{m}). It also agrees with its subject (\textsc{3} \textsc{sg} \textsc{f}) through the \textit{{}-o} suffix. In many languages we can specify items (verbs in this instance) simply; that is, for each feature there will be one value. But this will not work in Mian; the verb in \REF{ex:key:9} is both masculine and feminine. And just specifying both gender values is insufficient, we need to distinguish the two; it is masculine for the subject and feminine for the object. As pointed out by \citet{Zwicky1986}, accounting for inflection may require multiple feature specifications; in Mian we need separate feature specifications for subject agreement and for object agreement.\footnote{This does not imply that other featural specifications are unstructured. We may find, for instance, that gender values are distinguished only in the singular (as in Russian); and the Polish paradigm of \textit{muzeum} ‘museum’ in \tabref{tab:key:5} shows case values being distinguished only in the plural. But these are conditions imposed by feature \textit{values} (plural, in these instances); multiple feature values involve complete features. Thus in Mian, object agreement (of the type in \REF{ex:key:9}) requires person, number and gender (just as subject agreement does).} (Mian does not require an overt subject in examples like \REF{ex:key:9}.) We may find multiple specifications for inherent and contextual features (as in Upper Sorbian, \citealt{Corbett1987}) and for agreement with different controllers (as in \REF{ex:key:9}).\footnote{ \textrm{As a related phenomenon, potentially involving multiple feature specifications, in Nuer just a few nouns retain a construct state form, while all nouns (including those with the construct state form) inflect for case (Matthew Baerman, personal communication 15 \citealt{June2023}).}} Typically these distinct featural specifications arise from different syntactic requirements. 

We shall see that one set of forms can occur without the other, which helps confirm that these features are structured. While Mian verbs all have obligatory subject agreement, they may be uninflectable for object agreement, as in \REF{ex:key:10}:

Mian \citep[284]{Fedden2022}

\ea%10
    \label{ex:key:10}
    \gll \\
         \\
    \glt
    \z % you might need an extra \z if this is the last of several subexamples

          máam=e  bou-n-o=be

  mosquito=\textsc{sg.m}  swat-\textsc{real-3sg.f.sbj=decl}

  ‘She swatted the mosquito.’

More generally, the subject agreement forms (various specifications) appear on all verbs; the object agreement forms (various specifications) appear on some verbs (as in \REF{ex:key:9}) and not on others (as in \REF{ex:key:10}). Thus there is a slab of the paradigm which does not appear with certain verbs.

These data are sufficient to make the main point, but Mian is intriguingly more complex in having a second system of nominal classification. This has been called a system of verbal classifiers; however, the difference between gender and classifier is not clear-cut (Corbett, \citealt{FeddenFinkel2017}), and we shall call this system ‘\textsc{gender} \textsc{2}’. An example is given in \REF{ex:key:11}:

Mian \citep[288]{Fedden2022}

\ea%11
    \label{ex:key:11}
    \gll \\
         \\
    \glt
    \z % you might need an extra \z if this is the last of several subexamples

         máam=e  \textbf{dob}{}-ò-n-o=a

  mosquito=\textsc{sg.m}  \textbf{\textsc{3sg.gender2.obj}}{}-take.\textsc{pfv-seq-3sg.f.sbj=med}

  ‘She picked up the mosquito and then ...’

This second gender system is found mainly with verbs of object handling and movement: ‘give’, ‘take’, ‘put’, ‘throw’, ‘lift’, ‘turn’, ‘fall’ (some 50 verbs in all). For the verbs involved, gender 2 is obligatory. This system works on an absolutive basis, but only one verb is intransitive, and the remaining verbs show agreement with the object.\footnote{\textrm{There are also seven verbs which arguably show verbal number, alternating according to the number of the object \citep[311-314]{Fedden2019}.}} Like the object agreement discussed above, agreement is found prefixally. The form in \REF{ex:key:11} is the masculine (‘mosquito’ is masculine in this system). In brief, all verbs take subject agreement, and some additionally take object agreement. This object agreement is of two types: (i) some verbs take object agreement including gender as in \REF{ex:key:9}, while for some this slab of the paradigm is missing; (ii) some take object agreement including gender 2 as in \REF{ex:key:11}, and some do not. Finally, some verbs are uninflectable for both types of object agreement (see Corbett, \citealt{FeddenFinkel2017} for the detail). The canonically complete verb, with subject agreement and both types of object agreement, does not exist in Mian. There are no external conditions (\sectref{sec:key:5}) on whether a verb inflects for the object, and this has no effect on syntax (\sectref{sec:key:6}). 

Let us turn to the number of verbs involved (as shown in \tabref{tab:key:7}):


\begin{tabularx}{\textwidth}{XXX}

\lsptoprule

Transitive verb type (object agreement) & n & percent\\
never agree, e.g. \textit{bou} ‘swat’, as in \REF{ex:key:10} & \raggedleft 248 & \raggedleft 84\%\\
agree in person, number and gender, e.g. \textit{nâ’} ‘hit, kill’, as in \REF{ex:key:9} & \raggedleft 7 & \raggedleft 2\%\\
agree in gender 2, e.g. \textit{ò} ‘take’, as in \REF{ex:key:11} & \raggedleft 40 & \raggedleft 14\%\\
\lspbottomrule
\end{tabularx}
%%please move \begin{table} just above \begin{tabular
\begin{table}
\caption{Object agreement in Mian verbs: types \citep[293]{Fedden2022}}
\label{tab:key:7}
\end{table}

We see from \tabref{tab:key:7} that the number of verbs inflecting for object agreement is low. The key point, however, is that all verbs inflect for subject agreement, and a minority inflect for object agreement (in different ways). We return to Mian in \sectref{sec:key:6.1.2}. 

Before leaving multiple feature specifications in a slab of the paradigm, we note three further examples. The first is object agreement in Nankina (nnk), a Trans New Guinea language. Nine verbs have a full paradigm of object agreement, some others distinguish number only, and the rest are uninflectable for object agreement (\citealt{SpauldingSpaulding1994}: 40-42, \citealt{Corbett2012}: 164n15, \citealt{FeddenPaciaroni2022}). The second is possessed nouns in Dano (aso), also a Trans New Guinea language. Some of possessed nouns inflect for possessor person and number, and for case, while others inflect for possessor person and number, but not case (\citealt{Strange1972}, Baerman, \citealt{BrownCorbett2017}: 10-12, 28). Additional examples can be found in \citet{Windschuttel2019}. The third is number and gender agreement in Archi; \tabref{tab:key:4} showed that in different parts of speech, some items show agreement while others do not. This cross-cuts other inflectional possibilities, so that some items are fully uninflectable and some are only partially uninflectable; for present purposes, the number and gender agreement realizes one feature specification, and there are others which may or may not be realized.

We have established two dimensions of partial inflectability. First splits in the paradigm, where we contrasted slab and cells. Second the distinction in the system of feature specification, between a simple featural specification and multiple feature specification. We have seen that Mian shows multiple feature specification involving a slab of the paradigm. The remaining possibility would be multiple feature specification involving the cell(s) in the paradigm. This configuration seems highly unlikely, and yet it does occur. For this we return to Archi. Consider example \REF{ex:key:12}: 

Archi (\citealt{Kibrik1994}: 349, repeated from \REF{ex:key:5})

\ea%12
    \label{ex:key:12}
    \gll \\
         \\
    \glt
    \z % you might need an extra \z if this is the last of several subexamples

          buwa-mu  \textbf{b}{}-ez  ditːa<b>u  χːʷalli  a<b>u

  mother(\textsc{ii})-\textsc{sg.erg}  \textbf{\textsc{sg.iii}}{}-1\textsc{sg}.\textsc{dat}  early<\textsc{sg.iii}>  bread(\textsc{iii)[sg.abs]}  make<\textsc{sg.iii}>\textsc{pfv}

  ‘Mother made bread for me early.’

The pronoun \textit{bez} is first person singular, and it stands in the dative. Besides this specification for person, number and case, it has a further featural specification for the number and gender of the absolutive argument, with which it agrees. However, this is not general across the personal pronouns; rather it is limited to a few cells of the paradigm (see Table~8).


\begin{tabularx}{\textwidth}{XXXXXX}

\lsptoprule
%\hhline%%replace by cmidrule{~-----}
\multicolumn{1}{c}{~

} & \multicolumn{2}{c}{ \textsc{singular}} & \multicolumn{3}{c}{ \textsc{plural}}\\
& \textsc{1} \textsc{person}  & \textsc{2} \textsc{person}  & \multicolumn{2}{c}{ \textsc{1} \textsc{person}} & \textsc{2} \textsc{person} \\
&  &  & \multicolumn{1}{c}{\textsc{exclusive}} & \textsc{inclusive} & \\
\textsc{absolutive} & zon & un & \multicolumn{1}{c}{nen}  & nen<t’>u & žʷen \\
\textsc{ergative} & zari &  &  & nena<w>(u)

nena<r>u

nena<b>u

nen<t’>u  etc. & \\
\textsc{genitive}  & %%[Warning: Draw object ignored]
%%[Warning: Draw object ignored]
w-is      

d-is   

%%[Warning: Draw object ignored]
b-is      

is & wit & \multicolumn{1}{c}{ulu

d-olo

b-olo

olo      etc.}  & la<w>u

la<r>u

la<b>u

la<t’>u     etc. & wiš \\
\textsc{dative} & %%[Warning: Draw object ignored]
%%[Warning: Draw object ignored]
w-ez 

d-ez  

%%[Warning: Draw object ignored]
\textbf{b-ez}  

ez & wa-s & \multicolumn{1}{c}{w-el

d-el

b-el

el         etc.} & w-ela<w>(u)

d-ela<r>u

b-ela<b>u

el<t’>u     etc. & wež \\
\textsc{comitative} & za-ɬːu & wa-ɬːu & \multicolumn{2}{c}{la-ɬːu}  & žʷa-ɬːu \\
\textsc{similative}  & za-qˤdi & wa-qˤdi & \multicolumn{2}{c}{la-qˤdi}  & žʷa-qˤdi \\
\textsc{comparative}  & za-χur & wa-χur & \multicolumn{2}{c}{la-χur}  & žʷa-χur \\
\textsc{substitutive}  & za-kɬ’ena & wa-kɬ’ena & \multicolumn{2}{c}{la-kɬ’ena}  & žʷa-kɬ’ena \\
\textsc{superessive}  & za-t & wa-t & \multicolumn{2}{c}{la-t}  & žʷa-t \\
\textsc{superelative}  & za-tːi-š & wa-tːi-š & \multicolumn{2}{c}{la-tːi-š}  & žʷa-tːi-š \\
\textsc{superlative}  & za-tːi-k & wa-tːi-k & \multicolumn{2}{c}{la-tːi-k}  & žʷa-tːi-k \\
\textsc{superterminative}  & za-tːi-kǝna & wa-tːi-kǝna & \multicolumn{2}{c}{la-tːi-kǝna}  & žʷa-tːi-kǝna \\
\textsc{contelative} & za-ra-š & wa-ra-š & \multicolumn{2}{c}{la-ra-š}  & žʷa-ra-š \\
\textsc{contlative}  & za-ra-k & wa-ra-k & \multicolumn{2}{c}{la-ra-k}  & žʷa-ra-k \\
\textsc{contallative} & za-r-ši & wa-r-ši & \multicolumn{2}{c}{la-r-ši}  & žʷa-r-ši \\
\textsc{contterminative}  & za-ra-kǝna & wa-ra-kǝna & \multicolumn{2}{c}{la-ra-kǝna}  & žʷa-ra-kǝna \\
\lspbottomrule
\end{tabularx}
%%please move \begin{table} just above \begin{tabular
\begin{table}
\caption{Archi personal pronouns (partial paradigms), based on \citet[257-260]{Kibrik1977} and Chumakina \& \citet[102]{Corbett2015}}
\label{tab:key:8}
\end{table}

Note:   ‘etc.’ in a cell indicates the pattern of syncretism where the plural of \textsc{genders} \textsc{i} and \textsc{ii} matches \textsc{gender} \textsc{iii} \textsc{singular}, and the plural of \textsc{genders} \textsc{iii} and \textsc{iv} matches \textsc{gender} \textsc{iv} \textsc{singular.} 

  ‘\textsc{cont’} in the names of spatial cases indicates contact location.

The paradigms in \tabref{tab:key:8} are partial since some thirty case values are omitted. Recall the agreeing dative singular pronoun in example \REF{ex:key:10} above, indicated in bold in \tabref{tab:key:8}. This shows that certain cells in these paradigms inflect for agreement, while the majority do not.\footnote{ \textrm{For the remarkable first person plural inclusive absolutive case form in \tabref{tab:key:8}, which agrees with itself, see \citet[68-69]{Corbett2006}.}} To be clear what is going on, each of the cells shaded in yellow in the main paradigm houses a sub\nobreakdash-paradigm; for instance, the forms of the first person singular dative are structured as in Table~9:


\begin{tabularx}{\textwidth}{XXXX}

\lsptoprule
%\hhline%%replace by cmidrule{~~--}
\multicolumn{1}{c}{} &  & \multicolumn{2}{c}{ \textsc{number}}\\
%\hhline%%replace by cmidrule{~~--}
\multicolumn{1}{c}{} &  & \textsc{singular} & \textsc{plural}\\
 \textsc{gender} & \textsc{i} & w-ez & b-ez\\
& \textsc{ii} & d-ez & \\
& \textsc{iii} & b-ez & ez \\
& \textsc{iv} & ez & \\
\lspbottomrule
\end{tabularx}
%%please move \begin{table} just above \begin{tabular
\begin{table}
\caption{The first singular dative pronoun in Archi}
\label{tab:key:9}
\end{table}

Each of the shaded cells in \tabref{tab:key:8} includes an additional paradigm comparable to \tabref{tab:key:9}, with the same pattern of syncretism. (For comparison with other languages of the family see \citealt{KibrikKodzasov1990}: 220-223.) This is precisely the situation which seemed unlikely. Its existence means we have a full typology, as we see in \sectref{sec:key:4.3}.

\subsection{\textbf{The} \textbf{typology} \textbf{of} \textbf{partial} \textbf{inflectability}}

Putting together the analyses in \sectref{sec:key:4.1} and \sectref{sec:key:4.2} we arrive at the typology shown in \figref{fig:key:2}. 


\begin{tabularx}{\textwidth}{XXXX}

\lsptoprule
%\hhline%%replace by cmidrule{~~--} &  & \multicolumn{2}{c}{ definition of segments}\\
%\hhline%%replace by cmidrule{~~--} &  & slab (featural) & cell(s) (morphomic)\\
\multicolumn{1}{c}{{ system}

} & simple specification & Polish \textit{muzeum} ‘museum’ & Russian \textit{èxo} ‘echo’\\
& multiple feature specification & Mian object agreement & Archi pronouns\\
\lspbottomrule
\end{tabularx}
\begin{figure}
\caption{Partial inflectability: Two dimensions of comparison: summary}
\label{fig:key:2}
\end{figure}

We see that our typology of partial inflectability is complete (\figref{fig:key:2}). This opens the way for future investigation of the distribution of examples over the four types. The degrees of inflectability found in each type also deserve further work: the examples found to date typically have fewer lexical items inflecting for the full featural possibilities than those which are partially uninflectable (see, for instance, \tabref{tab:key:7} on Mian verb types). 

\section{\textbf{Conditions} \textbf{(“incoming”} \textbf{factors)}}

In the canonical world, inflecting items realize their feature specifications fully, irrespective of external conditions. Thus a verb may have feature values for tense and aspect, and agreement specifications, and this is sufficient to determine its form. That is, external conditions such as word order are irrelevant: a verb does not change its tense value when in clause-initial position, for instance. However, in the real world we do find significant instances where items inflect or do not inflect, subject to external conditions (in addition to the expected featural requirements imposed, for instance, by government and agreement). Such conditions are non-canonical, since the syntax would be simpler without them. The result of such conditions has been termed ‘constructional non-inflectedness’ by \citet[144]{Spencer2020}, and ‘syntagmatic/contextual uninflectedness’ by Doleschal (2000-2001: 14). I first present the different types of conditions which lead to uninflectedness (full or partial), then in \sectref{sec:key:5.5} I review the forms which result. 

\subsection{\textbf{Syntactic} \textbf{position:} \textbf{German}}

A classic example is provided by German adjectives. In attributive position these inflect for case, number and gender, with two levels of reduced agreement \citep[95-96]{Corbett2006} to which we return in \sectref{sec:key:5.5}. Elsewhere, notably in predicate position, adjectives are uninflected:

\ea%13
    \label{ex:key:13}
    \gll \\
         \\
    \glt
    \z % you might need an extra \z if this is the last of several subexamples

          Dies-er  Artikel  ist  sehr  interessant 

  this\textsc{{}-sg.m.nom}  article(\textsc{m)[sg.nom]}  be.\textsc{prs}.\textsc{3sg}  very  interesting

‘this article is very interesting’

In \REF{ex:key:13} the adjective \textit{interessant} ‘interesting’ lacks any inflection.\footnote{ \textrm{For the complex issue of the conditions on case marking and lack of inflection for case in German, see \citet{Spencer2009} and references there.} } Hence the condition determining whether the adjective inflects is a syntactic one. 

\subsection{\textbf{Position} \textbf{in} \textbf{the} \textbf{phrase:} \textbf{Tsova-Tush}}

In several languages we find inflection restricted to the right edge of a nominal phrase (this may involve conjuncts or an attributive modifier and a noun). We might anticipate a simple external condition on inflection: the item on the right edge inflects, others do not. For instance, an attributive occurring phrase-finally would inflect, while if standing before a noun (hence not at the right edge), it would not. With this in mind, consider the data in \tabref{tab:key:10}, from the Nakh language Tsova-Tush (also known as Batsbi), as analysed by \citet{Sims-WilliamsBaerman2021}. 


\begin{tabularx}{\textwidth}{XXX}

\lsptoprule
%\hhline%%replace by cmidrule{~--}
\multicolumn{1}{c}{} & Substantivized adjective in final position ‘bad one’ & Attributive adjective in phrase ‘bad child’\\
\textsc{absolutive} & mos:i\textsuperscript{n} & mos:i\textsuperscript{n}   bader\\
\textsc{ergative} & muis:čo-v & muis:čO  badre-v\\
\textsc{genitive} & muis:čo-\textsuperscript{n} & muis:čO  badre-\textsuperscript{n}\\
\textsc{dative} & muis:čo-n & muis:čO   badre-n\\
\lspbottomrule
\end{tabularx}
%%please move \begin{table} just above \begin{tabular
\begin{table}
\caption{ Tsova-Tush (bbl, \citealt{HoliskyGagua1994}: 171-172, \citealt{Sims-WilliamsBaerman2021}: 28)}
\label{tab:key:10}
\end{table}

Note: \textsuperscript{n} indicates nasalization, O indicates a reduced vowel (word-final position)

In \tabref{tab:key:10} we see that the substantivized adjective in final position is indeed inflected; but when there is attributive adjective, followed by a noun, it is the noun which takes the inflection. Sims-William \& Baerman stress the word-order effect, given that group inflection is an area feature shared by various languages of the Caucaus. From this perspective, we have an external condition on the paradigm: inflection is found at the right edge of the nominal phrase (and note that the absolutive singular has the bare stem, a significative absence, \sectref{sec:key:2.3}). However, this adjective is one that also has a stem alternation, contrasting absolutive and oblique (the remaining case values). Even when it does not bear the main inflection, it retains some inflection, that of the stem alternation.

\subsection{\textbf{Featural} \textbf{condition:} \textbf{Sanskrit}}

We turn to a different type of condition, found in the Sanskrit (san) future, as presented by \citet{Stump2013}. The tense in question, termed the periphrastic future,\footnote{See further \citet{Lowe2017} for critical discussion of the periphrastic status of this future tense at different periods of Sanskrit and for further examples.}  consists of an agent noun and, in the first and second persons, a form of the copula; this is set out in \tabref{tab:key:11}.


\begin{tabularx}{\textwidth}{XXXX}

\lsptoprule
%\hhline%%replace by cmidrule{~~--} &  & \multicolumn{2}{c}{ components}\\
%\hhline%%replace by cmidrule{~~--} &  & verb derivative & copula\\
\multicolumn{1}{c}{ \textsc{person}} & 1 and \textsc{2}  & agent-noun derivative

\textsc{masculine}

\textsc{nominative}

\textsc{singular} & \textsc{present} \textsc{indicative} \textsc{active} (very rarely \textsc{middle}) forms of \textit{as} ‘be’\\
& 3 & agent-noun derivative 

\textsc{masculine}

\textsc{nominative}

\textsc{singular} \textsc{/} \textsc{dual} \textsc{/} \textsc{plural} & [none]\\
\lspbottomrule
\end{tabularx}
%%please move \begin{table} just above \begin{tabular
\begin{table}
\caption{Periphrastic future-tense formation in Sanskrit \citep[111]{Stump2013}}
\label{tab:key:11}
\end{table}

The forms of \textit{as} ‘be’ are given in \tabref{tab:key:12}. This verb has all forms available, though the third person forms do not appear in the periphrastic tense in question. 


\begin{tabularx}{\textwidth}{XXXX}

\lsptoprule
%\hhline%%replace by cmidrule{~---}
\multicolumn{1}{c}{} & \textsc{singular} & \textsc{dual} & \textsc{plural}\\
1 & ásmi & svás & smás\\
2 & ási & sthás & sthá\\
3 & ásti & stás & sánti\\
\lspbottomrule
\end{tabularx}
%%please move \begin{table} just above \begin{tabular
\begin{table}
\caption{ Present indicative active of \textit{as} ‘be’ in Sanskrit \citep[111]{Stump2013}}
\label{tab:key:12}
\end{table}

More importantly, for our purposes, the agent noun also has a full paradigm, distinguishing number, gender and case. With that in mind, consider the paradigm in \tabref{tab:key:13}.


\begin{tabularx}{\textwidth}{XXXX}

\lsptoprule
%\hhline%%replace by cmidrule{~---}
\multicolumn{1}{c}{} & \textsc{singular} & \textsc{dual} & \textsc{plural}\\
1 & dāt\'{ā} asmi & dāt\'{ā} svas & dāt\'{ā} smas\\
2 & dāt\'{ā} asi & dāt\'{ā} sthas & dāt\'{ā} stha\\
3 & dāt\'{ā} & \textbf{dāt\'{ā}rau} & \textbf{dāt\'{ā}ras}\\
\lspbottomrule
\end{tabularx}
%%please move \begin{table} just above \begin{tabular
\begin{table}
\caption{Periphrastic future-tense forms of \textit{dā} ‘give’, without sandhi (Stump: 2013: 111)}
\label{tab:key:13}
\end{table}

Presenting the paradigm in this form is helpful, since it shows the partial inflectability which interests us. Of the full paradigm of the agent nominal, we find just the number distinction, and that only for third person forms. This can be seen, though less obviously, in the ‘real’ forms, those which arise as a result of automatic sandhi processes, where the two parts are pronounced as a single phonological word (see \tabref{tab:key:14}). In the examples, sandhi is indicated by ‿.


\begin{tabularx}{\textwidth}{XXXX}

\lsptoprule
%\hhline%%replace by cmidrule{~---}
\multicolumn{1}{c}{} & \textsc{singular} & \textsc{dual} & \textsc{plural}\\
1 & dāt\'{ā}smi & dāt\'{ā}svas & dāt\'{ā}smas\\
2 & dāt\'{ā}si & dāt\'{ā}sthas & dāt\'{ā}stha\\
3 & dāt\'{ā} & \textbf{dāt\'{ā}rau} & \textbf{dāt\'{ā}ras}\\
\lspbottomrule
\end{tabularx}
%%please move \begin{table} just above \begin{tabular
\begin{table}
\caption{Periphrastic future-tense forms of \textit{dā} ‘give’, with sandhi \citep[111]{Stump2013}}
\label{tab:key:14}
\end{table}

The key point for our analysis is that the agent nominal has a full paradigm, but when used in this future tense it inflects partially: it does not inflect when the periphrastic verb is first or second person; for the third person the agent nominal inflects for number but not for gender (only the masculine is found).\footnote{ \textrm{We see distributed exponence here: where the agent nominal is uninflectable for number, the auxiliary is present, and inflects for number; when the auxiliary is absent (in the third person) the agent nominal marks number.} } Let us see how this works in examples:

Sanskrit \citep[112]{Stump2013}

\ea%14
    \label{ex:key:14}
    \gll \\
         \\
    \glt
    \z % you might need an extra \z if this is the last of several subexamples

          tatas  tvā  pārayitā‿asmi

  from.that  \textsc{2sg.acc}   rescue.\textsc{nmlz.sg.m}{}-be.\textsc{prs.1sg}

  ‘I will rescue (lit: I am rescuer) you from that.’   (\textit{Śatapatha} \textit{Brāhmaṇa} 1.8.1.2)

In \REF{ex:key:14} we have a first person singular subject, and the agent nominal is masculine singular. There is also a direct object (\textit{tvā} ‘you’), in the accusative, showing that this tense is not an ordinary predicate nominal construction. Compare the plural in \REF{ex:key:15}: 

Sanskrit (part of example in \citealt{Stump2013}: 113)

\ea%15
    \label{ex:key:15}
    \gll \\
         \\
    \glt
    \z % you might need an extra \z if this is the last of several subexamples

          mahac  ‿choka-bhayam  prāptā‿smaḥ\footnote{\textrm{\textit{smaḥ}} \textrm{is the pre-pausal sandhi variant of the first person plural} \textrm{\textit{smas} }\textrm{(Greg Stump, personal communication, 5 \citealt{August2013})}}

  great.\textsc{sg.acc}  pain-fear.\textsc{sg.acc}  meet.\textsc{nmlz.sg.m}{}-be.\textsc{prs.1pl}

  ‘We are going to meet with great pain and dread.’  (\textit{Gopatha Brāhmaṇa} 1.1.28)

Note that although we have a plural subject, the agent nominal remains singular here (since the verb is first person). 

Sanskrit \citep[113]{Stump2013}

\ea%16
    \label{ex:key:16}
    \gll \\
         \\
    \glt
    \z % you might need an extra \z if this is the last of several subexamples

          ekā  jananayitā  putram

  one.\textsc{sg.nom}.\textsc{f}  produce.\textsc{nmlz.sg.}\textbf{\textsc{m}}  son.\textsc{sg.acc}

  ‘One (feminine) will bear a son.’    (\textit{Rāmāyaṇa} 1.38.8)

Example \REF{ex:key:16} is significant: the subject is third person (so no copula); it is clearly feminine, but the agent nominal maintains masculine gender. However, the agent nominal can inflect (for number) in future tense use, as example \REF{ex:key:17} shows: 

Sanskrit \citep[114]{Stump2013}

\ea%17
    \label{ex:key:17}
    \gll \\
         \\
    \glt
    \z % you might need an extra \z if this is the last of several subexamples

          na  jāne  yātāras  tava  ripavaḥ  kena  ca  pathā.

  not  know.\textsc{1sg}  go.\textsc{nmlz.pl.m}  thy  enemy.\textsc{pl.nom}  which.\textsc{sg.ins}  and  path.\textsc{sg.ins}

  ‘I don’t know by which path your enemies will go.’   (\textit{Bhojaprabandha} 55)

Here we see a third plural subject and the agent nominal is plural. This fits with the pattern in \tabref{tab:key:11}. The agent nominal has a full paradigm, but when used for the future it does not inflect at all when the subject is first or second person; when the subject is third person it inflects for number \REF{ex:key:17} but not gender \REF{ex:key:16}. The agent nominal’s pattern of partial inflection is conditioned by the person of the subject, that is, by a featural specification. 

\subsection{\textbf{Further} \textbf{instances}}

Other examples of external conditions include the following varied selection. Proper names may be accompanied by classifier-like nouns (sortals) and the proper name may then remain uninflected. See Matushansky (2022, 2023) for careful discussion of these constructions in Russian. An example is given in \REF{ex:key:18}:

Russian \citep[235]{Matushansky2023}

\ea%18
    \label{ex:key:18}
    \gll \\
         \\
    \glt
    \z % you might need an extra \z if this is the last of several subexamples

          na  ulic-e /   *ulic-a  Jakimank-e /   Jakimank-a 

  on  street-\textsc{sg.loc} /   *street-\textsc{sg.nom}  Yakimanka-\textsc{sg.loc} /  Yakimanka-\textsc{sg.nom} 

‘on Yakimanka Street’

In this instance (see \citealt{Matushansky2023} for the other types), while the sortal noun must take the case value required by the syntactic context, the proper name may be in that case or may stand in the nominative. For appositives in Polish see \citet{Cetnarowska2021}.

In Kayardild, the external condition is that the set of words on which a morphosyntactic feature is realized is determined by concentric domains of features \citep[78-84]{Round2013}. In Lithuanian (\citealt{Arkadiev2013,Arkadiev2020}), participles agree or lack agreement morphology depending on the construction (and there is a complex web of interrelated constructions). \citet{Giger2016} describes the loss of inflection of participles as a participial perfect emerges in Slovak. And for the interesting complexities of Hungarian inflecting infinitives and Latvian genitive compounds see \citet{Spencer2025}.

\subsection{\textbf{Forms} \textbf{used} \textbf{and} \textbf{distinctions} \textbf{retained} \textbf{in} \textbf{constructional} \textbf{non-inflectedness}}

We have examined the types of construction which condition (often partial) non\nobreakdash-inflectedness. We now switch our focus to the forms used in examples of constructional non\nobreakdash-inflectedness, and the distinctions retained. The clearest type is that found in the German predicate adjective (example \REF{ex:key:13}); this takes a form lacking all inflection, and retaining none of the featural distinctions of inflected forms. This implies two distinctions. First, there are uniquely uninflected forms, as in our German example, opposed to normal inflected forms, used for constructional non\nobreakdash-inflectedness, such as the Russian nominative (\sectref{sec:key:5.4}). Second, we contrast forms where no distinctions are retained, as in the German example, opposed to those where at least one distinction is retained, as in Tsova-Tush. Significant examples are given in \tabref{tab:key:15}, and then key points are highlighted.


\begin{tabularx}{\textwidth}{XXXX}

\lsptoprule

example & unique vs 

normal forms & featural distinctions retained & source\\
German predicative adjective & unique & none & see \sectref{sec:key:5.1}\\
Tsova-Tush attributive adjective (not phrase-final) & unique & \textsc{absolutive} vs \textsc{oblique} & see \sectref{sec:key:5.2}\\
Somali verb (subject focus) & unique & \textsc{3sg.f} \textsc{/} \textsc{1pl} / rest & Corbett (2006: 93, 203-204) and references there \\
Inari Saami reduced agreement & normal & \textsc{sg} vs \textsc{du-pl} & \citet[94-95]{Corbett2006}, \citet{Toivonen2007}\\
Russian sortal noun & normal: \textsc{nominative} & \textsc{sg} vs \textsc{pl} & see \sectref{sec:key:5.4}\\
Sanskrit periphrastic future & normal: \textsc{masculine} (and \textsc{singular} in person 1 and 2) & \textsc{sg} vs \textsc{du} vs \textsc{pl} (in person 3) & see \sectref{sec:key:5.3}, \citet{Stump2013}\\
German attributive adjective & normal & \textbf{\textsc{case}} and \textbf{\textsc{number}} (limited degree) & \citet[95-96]{Corbett2006}\\
\lspbottomrule
\end{tabularx}
%%please move \begin{table} just above \begin{tabular
\begin{table}
\caption{Forms and distinctions in constructional non\nobreakdash-inflectedness}
\label{tab:key:15}
\end{table}

Somali deserves mention since it has reduced agreement forms which are unique to the reduced paradigm; furthermore, the distinctions made in the reduced paradigm are not a simple collapsing on those in the full paradigm (as in Tsova-Tush), but rather the patterns of syncretism are different. Inari Saami is included because the reduced paradigm has no unique forms; it retains the third person singular and the third plural from the full paradigm, and these fill out the reduced paradigm. Reduced agreement (\citealt{FeddenPaciaroni2022}, and \citealt{Loporcaro2025}), a subset of partial inflection determined by external conditions, is a fruitful area to work on, because agreement allows us to manipulate the feature values involved, which can be useful for establishing the nature of the partially inflected forms. Continuing with the data in \tabref{tab:key:15}, the Sanskrit periphrastic future (\sectref{sec:key:5.3}) is illuminating because we find possible inflected forms, but used beyond their normal range. Here we can ask whether we have a default form (an inflected form used when there is no specification). This is a plausible analysis for the Sanskrit future, where we find masculine as the gender default, and singular as the number default. (The nominative is expected as the case value in the predicate.) However, in the third person, distinctions of number are retained: we have a reduced paradigm rather than a single default form. Thus the features in Sanskrit have sufficient values for it to be clear what is going on; we need to beware of small systems, which can be misleading.\footnote{ \textrm{There is careful discussion of default agreement vs lack of agreement inflection in Lithuanian in Arkadiev (2020: 383-387, 403-404, 408-412). And we need to establish the default instance by instance; thus in Serbo-Croat (\citealt{Browne1993a}: 329, 374) and Polish the accusative is an attractor for length-of-time and quantity constructions.}} Finally, we need to note the complexities of German adjectives in attributive position. There are two stages of reduction: the paradigms are traditionally labelled strong, mixed and weak. There is a notional paradigm with 16 cells (four case values, plural vs singular and – in the singular – three gender values). However, there are extensive syncretisms so that even the strong paradigm has only five distinct forms. The mixed paradigm has four of those forms, and the weak has only two; however, we do not find a straightforward collapsing of distinctions, but rather we see different patterns of syncretism. The mixed and weak paradigms retain some distinction within each of the features (case, number and person).

We have observed that there are lexemes which do not inflect, even though this might be expected (\sectref{sec:key:3}), and this uninflectability is a matter of degree (\sectref{sec:key:4}). In this section (\sectref{sec:key:5}), we have seen that there are external conditions which determine whether otherwise inflecting items inflect. And here too it is a matter of degree. We now turn to the final logical possibility, the possible outward effect of uninflectability. 

\section{\textbf{External} \textbf{consequences} \textbf{(“outgoing”)}}

Since uninflectability is a matter of inflection, we would expect no syntactic consequences. And we do find instances of this canonical situation (\sectref{sec:key:6.1}). But there are also claims in the literature that uninflectability has consequences, and I discuss some of these (\sectref{sec:key:6.2}).

\subsection{\textbf{No} \textbf{external} \textbf{consequences}}

Often this situation is simply assumed: researchers report the inflectional situation and imply that there is no more to be said. However, there are also studies which look carefully for external consequences, and do not find any. We consider in turn \citet{Nichols2018} on Ingush, and \citet{Fedden2022} on Mian.

\subsubsection{\textbf{Agreement} \textbf{in} \textbf{Ingush} \textbf{(Nakh-Dagestanian)}}

Ingush (inh), which has some 300,000 speakers living mainly in Ingushetia in the North Caucasus, is a Nakh language (hence a relative of Tsova-Tush, discussed in \sectref{sec:key:5.2} and a more distant relative of Archi, introduced in \sectref{sec:key:1.3} and \sectref{sec:key:3.1.2}). It has ‘an arbitrary and strictly lexical bifurcation of verbs’ \citep[845]{Nichols2018}, in that some have gender and number agreement and some do not. There are 400 simple verbs at most, of which around 30\% show gender and number agreement (change of the initial consonant); however, there are other possibilities, including light verbs and tense auxiliaries, so that inflected verb forms have some gender and number agreement in probably well over 50\% of the instances (\citealt{Nichols2011}: 141n63, 235).\footnote{ \textrm{That is 90/294 (30.6\%) of the simple verb roots; if those with number distinctions (verbal number) are included, the figure is 109/393 (27.7\%). Of the simple adjectives, fewer than 10\% take agreement \citep[141]{Nichols2011}. See \citet[141-145]{Nichols2011} for details on the relation of gender and number in agreement.} } Given this situation, we can use variation in one phenomenon (inflection for agreement or lack of it) to test proposed explanations for another. Thus it is sometimes suggested that the function of agreement is to facilitate reference tracking. By comparing agreeing with non-agreeing verbs, we can ask whether we can see differences in other means of reference tracking, in particular the presence of an overt argument (which might be introduced to facilitate reference tracking when the verb does not agree). \citet{Nichols2018} has measured carefully, using a corpus of narrative texts (mainly spoken). The results are summarized in \tabref{tab:key:16}.


\begin{tabularx}{\textwidth}{XXXXXX}

\lsptoprule
%\hhline%%replace by cmidrule{~~----}
\multicolumn{1}{c}{~} & ~ & \textsc{s} & \textsc{a} & \textsc{o/t} & total\\
verb root agrees & n & \raggedleft 100 & \raggedleft 41 & \raggedleft 37 & \raggedleft 178\\
& \% overt argument & \raggedleft 88 & \raggedleft 59 & \raggedleft 92 & \raggedleft 82\\
 verb root does not agree & n & \raggedleft 41 & \raggedleft 38 & \raggedleft 24 & \raggedleft 103\\
& \% overt argument & \raggedleft 85 & \raggedleft 76 & \raggedleft 83 & \raggedleft 82\\
\lspbottomrule
\end{tabularx}
%%please move \begin{table} just above \begin{tabular
\begin{table}
\caption{Ingush: verb agreement and overt arguments (tokens) (from \citealt{Nichols2018}: 858)}
\label{tab:key:16}
\end{table}

Note:   S: single argument of intransitive verb; A: agent-like argument of a monotransitive; O: direct object of a monotransitive; T: theme-like or patient-like object of a ditransitive

We see that there is no difference between the examples with an agreeing verb and those with a non-agreeing verb. The morphological difference is not reflected in the syntax. For comparable results from the Dagestanian language Lak, see \citet{Forker2018}.

\subsubsection{\textbf{Agreement} \textbf{in} \textbf{Mian}}

Mian was introduced in \sectref{sec:key:4.2}. Based on a fully annotated Mian test corpus of 14 texts, some 4,000 words, with around 1,700 clauses, \citet{Fedden2022} undertook a study similar to that of \citet{Nichols2018}, discussed in \sectref{sec:key:6.1.1}. We are concerned only with object agreement, since all verbs show subject agreement. There are three types of transitive verb; recall from \tabref{tab:key:7} that only 16\% of Mian verbs show object agreement, 2\% showing agreement in person, number and gender and 14\% taking gender 2. Thus object agreement is relevant for a lower proportion of the transitive verbs than that found in Ingush. Due to the high type and token frequency of non-agreeing verbs, the syntax is constantly encountering gaps. What then can we learn about reference tracking? The data on the presence of overt arguments are given in Table~17.


\begin{tabularx}{\textwidth}{XXXXX}

\lsptoprule
%\hhline%%replace by cmidrule{~~---}
\multicolumn{1}{c}{~} & ~ & gender & gender 2 & combined\\
verb agrees~ & n & \raggedleft 61 & \raggedleft 218 & \raggedleft 279\\
& \% overt & \raggedleft 33 & \raggedleft 39 & \raggedleft 38\\
verb does not agree & n &  &  & \raggedleft 397\\
& \% overt &  &  & \raggedleft 43\\
\lspbottomrule
\end{tabularx}
%%please move \begin{table} just above \begin{tabular
\begin{table}
\caption{Mian object agreement: verb agrees or not \citep[300]{Fedden2022}}
\label{tab:key:17}
\end{table}

The difference between the highlighted cells is not statistically significant. And so, although the situation in Mian is rather different from that in Ingush, we find again that inflectability has no syntactic effect.

\subsection{\textbf{External} \textbf{(syntactic)} \textbf{consequences}}

There are various places in the literature where it is claimed that uninflectability does have syntactic consequences. I will report the data on some of the suggested examples (I have carried out only limited fieldwork and cannot vouch for them). I will note the others briefly. Typically there is speaker variation. 

\subsubsection{\textbf{Serbo-Croat} \textbf{\textit{doba}} \textbf{‘time,} \textbf{season,} \textbf{age,} \textbf{period’}}

This noun is unusual in several ways, including the fact that it is the only neuter noun in the language with the nominative singular in \textit{{}-a.} In a language where uninflectability is rare (\citealt{Browne1993a}: 378, 1993b), this noun is often cited in grammars and dictionaries as uninflectable, with \textit{doba} as its unique form. However, the situation is more complex, as was pointed out as early as \citet[246]{Kovačević1951}, see also \citet[212-214]{Stevanović1975} and Nikolić (1995-1996: 16-17). The potential paradigm is given in \tabref{tab:key:18}, though the use of the form \textit{doba} is found in other cells of the paradigm too.


\begin{tabularx}{\textwidth}{XXX}

\lsptoprule
%\hhline%%replace by cmidrule{~--}
\multicolumn{1}{c}{} & \textsc{singular} & \textsc{plural}\\
\textsc{nominative} & doba & doba\\
\textsc{accusative} & doba & doba\\
\textsc{genitive} & doba & dobā\\
\textsc{dative} & dobu & dobima\\
\textsc{locative} & dobu & dobima\\
\textsc{instrumental} & dobom & dobima\\
\lspbottomrule
\end{tabularx}
%%please move \begin{table} just above \begin{tabular
\begin{table}
\caption{Inflection of Serbo-Croat \textit{doba} ‘time’ \citep[79]{Benson1971}}
\label{tab:key:18}
\end{table}

Our concern here is those instances when the noun remains uninflectable, since attributive modifiers can be affected by the uninflectability of the noun. There are examples in both \citet{Kovačević1951} and \citet{Stevanović1975} of phrases with \textit{doba} ‘time’, governed by a preposition, where the attributive modifier appears to follow the form of \textit{doba} rather than standing in the appropriate case. Consider these alternatives:

Serbo-Croat (\citealt{Kovačević1951} and web sources)

\ea%19
    \label{ex:key:19}
    \gll \\
         \\
    \glt
    \z % you might need an extra \z if this is the last of several subexamples

        a  u  t-om  doba  \REF{ex:key:19}b  u  t-o  doba

  in  that-\textsc{sg.loc.n/m}  time    in  that\textsc{{}-sg.acc.n} time

  ‘at that time’  ‘at that time’

The use of the locative modifier (with the uninflected \textit{doba}) is rare \REF{ex:key:19a}; much more common is the use of the \textit{to} form of the demonstrative \REF{ex:key:19b}. We can treat this as the accusative, and note two factors here. First the accusative tends to be used, rather than the locative, more generally in such time expressions: this is true both for other attributive modifiers of \textit{doba} ‘time’ \citep[11-13]{Browne2015} and for other nouns denoting time which do not have the unusual characteristics of \textit{doba}. And second, \textit{to} \textit{doba} ‘that time’ could be analysed as moving towards being an adverbial expression, a fixed phrase unaffected by external government. We therefore take a preposition governing the genitive case: then the first factor no longer applies, while the second does. \citet{Browne2015} reports a web count of three relevant examples:

\citet[12]{Browne2015}: web count 31.12.2014

\ea%20
    \label{ex:key:20}
    \gll \\
         \\
    \glt
    \z % you might need an extra \z if this is the last of several subexamples

          od  t-og  doba    285 instances: 43.5\%

  from  that-\textsc{sg.gen.n/m}  time

  ‘from that time’

\ea%21
    \label{ex:key:21}
    \gll \\
         \\
    \glt
    \z % you might need an extra \z if this is the last of several subexamples

          od  t-oga   doba    242 instances: 36.9\%

  from  that-\textsc{long}.\textsc{sg.gen.n/m}  time

  ‘from that time’

\ea%22
    \label{ex:key:22}
    \gll \\
         \\
    \glt
    \z % you might need an extra \z if this is the last of several subexamples

          od  t-o  doba    128 instances: 19.5\%

  from  that-\textsc{sg.nom/acc.n}  time

  ‘from that time’

Since \textit{od} ‘from’ takes the genitive, \REF{ex:key:20} and \REF{ex:key:21} are as expected. But a minority of instances as in \REF{ex:key:22} show that \textit{doba} ‘time’ does not always behave as we would expect from an uninflectable noun. Rather we find a syncretic nominative-accusative form, which we take as accusative for the reasons given above. But this would imply that we have different behaviours, depending on the particular cell of the paradigm (\sectref{sec:key:4.1}). Given that, the suggestion that \textit{to doba} ‘that time’ is becoming a fixed adverbial expression of time, impervious to case government, seems a reasonable one. 

\subsubsection{\textbf{Serbo-Croat} \textbf{numeral} \textbf{phrases}}

Serbo-Croat numerals have largely lost inflection. The higher numerals, \textit{pet} ‘five’ and above, typically do not inflect (and they take the genitive plural of the quantified noun).\footnote{ \textrm{They lost inflection in the 16}\textrm{\textsuperscript{th}} \textrm{to 17}\textrm{\textsuperscript{th}} \textrm{centuries, according to \citet[138]{Rogić1954}. However, the higher, more noun-like numerals (}\textrm{\textit{stotina}} \textrm{‘hundred’,} \textrm{\textit{tisuća}} \textrm{and} \textrm{\textit{hiljada}} \textrm{‘thousand’,} \textrm{\textit{milion} }\textrm{and} \textrm{\textit{milijun}} \textrm{‘million’) can inflect, particularly in complex numerals, or stand in the accusative \citep[11-12]{Popović1979}.} } The numerals \textit{dva} ‘two’ (masculine and neuter with \textit{dve} feminine), \textit{tri} ‘three’ and \textit{četiri} ‘four’ are losing inflection for case, but have not done so completely, and this has engendered a substantial literature.\footnote{\citet{Marković1951} pointed out that \textit{dva} ‘two’, and even more so, \textit{tri} ‘three’ and \textit{četiri} ‘four’ were becoming uninflectable for case; this change was well under way in the 19\textsuperscript{th} century, and he discusses instances where inflection is preserved in relevant texts. He noted already the tendency for numeral phrases which would be in the instrumental to be replaced with \textit{s} ‘with’ and the uninflectable form of the numeral. \citet{Rogić1954} gives examples illustrating the loss of inflection of these numerals. \citet{Hamm1956} discusses grammarians’ statements on the topic and usage in the press. M. \citet{Popović1967} talks of ‘blocks’ which inflect, meaning that in constructions consisting of preposition + numeral + noun, the numeral and noun form a ‘block’ whose case is registered in the preposition. The journal \textit{Naš jezik} had three papers on the topic in 1976: \citet{Kraljević1976} gives a wealth of data, including dialect data, showing that lack of inflection is more prevalent in the east; \citet{Nikolić1976} gives summary information on 28 dialects, showing that inflection of the numeral is more likely if the noun is feminine, more likely in the dative, locative and instrumental cases than in the genitive, more likely with \textit{dva} ‘two’ than with \textit{tri} ‘three’, and more likely with \textit{tri} ‘three’ than with \textit{četiri} ‘four’; \citet{Radović-Tešić1976} concentrates on the range of forms of the numerals. Lj. \citet{Popović1979} discusses uninflectability within the general morphosyntactic system of the numerals; his papers on uninflectability and the dative are discussed in the main text. Ďurovič (1985) also gives a general account, including uninflectedness, as does \citet{Browne2022}, while \citet{Grubišić1994} discusses the grammatical tradition on the topic. \citet{Kuzmić2007} examines Čakavian legal texts (13\textsuperscript{th}{}-18\textsuperscript{th} centuries) while \citet{KuzmićKuzmić2007} analyse Kajkavian legal texts (16\textsuperscript{th}{}-18\textsuperscript{th} centuries); both papers provide a wealth of textual examples from the respective varieties, including evidence on the rise of uninflectability. \citet[148-193]{Stefanović2019} provides many examples, of both inflectedness and uninflectedness.} These lower numerals require the ‘adnumerative’ of the quantified noun, and this is realized as the genitive singular or the nominative plural, depending on the inflection class of the noun \citep[159-160]{Corbett2009}. For our purposes, the typical situation is the interesting one, that where numeral phrases increasingly function as non-inflecting units; that is, the numeral does not inflect, and it determines the inflectional form of the quantified noun, giving a constituent which is impervious to external syntax. In other words, a numeral + noun unit (given in square brackets in our examples) behaves somewhat like an uninflectable noun. Such phrases are fine when they occur in syntactic positions requiring the nominative, accusative or genitive case; the latter is illustrated in \REF{ex:key:23}:

Serbo-Croat (\citealt{WechslerZlatić2003}: 150)

\ea%23
    \label{ex:key:23}
    \gll \\
         \\
    \glt
    \z % you might need an extra \z if this is the last of several subexamples

          seć-am  se  [ova četiri glumca]\footnote{ \textrm{Here} \textrm{\textit{glumca}} \textrm{‘actor’ is in the adnumerative, realized as the genitive singular;} \textrm{\textit{ova}} \textrm{‘these’ is also in the adnumerative, realized as the nominative plural neuter \citep[17-18]{Corbett2023}.}}    [\textsc{genitive}]

  remember-\textsc{prs}.1\textsc{sg}  \textsc{refl}  [\textsc{dem} four actor]

  ‘I remember these four actors.’

They are also fully acceptable with prepositions taking various cases, as in \REF{ex:key:24}

Serbo-Croat (\citealt{WechslerZlatić2003}: 150)

\ea%24
    \label{ex:key:24}
    \gll \\
         \\
    \glt
    \z % you might need an extra \z if this is the last of several subexamples

          razgovar-am   sa  [ova četiri glumca]    [\textit{sa} + \textsc{instrumental}]

  talk-\textsc{prs}.1\textsc{sg}  with  [\textsc{dem} four actor]

  ‘I am talking to these four actors’

However, these numeral phrases do not fit everywhere. As pointed out by Lj. Popović (1981-1982, see also 1982), the case frame matters: they do not work as a dative:\footnote{\textrm{I give examples from the extended discussion by Wechsler \& Zlatić, since they are simpler than those in the original. See also \citet[374-375]{Browne1993a}. As for the remaining two case values: (i) the locative is found only with prepositions, and so it causes no problem; (ii) the instrumental, interestingly, can be replaced in these constructions by} \textrm{\textit{sa} }\textrm{‘with’ plus the uninflected unit, to give an acceptable form comparable to \REF{ex:key:24}. There is no comparable escape for the dative (Lj. \citealt{Popović1981}-1982: 606).}}

Serbo-Croat (\citealt{WechslerZlatić2003}: 150)

\ea%25
    \label{ex:key:25}
    \gll \\
         \\
    \glt
    \z % you might need an extra \z if this is the last of several subexamples

          *dati  knjig-e  [ova četiri studenta]    [*\textsc{dative}]

  give.\textsc{inf}  book-\textsc{pl.acc}  [\textsc{dem} four student]

  ‘to give books to these four students’

This is indeed surprising; in some case frames the lack of morphological realization of the case value is unproblematic, while for the dative there is an issue \REF{ex:key:25}. The non-inflecting unit behaves as though it is defective, lacking a dative case form. There is further discussion in Wechsler \& \citet[148-152]{Zlatić2003}, Franks (forthcoming) and for higher numerals in \citet{Bošković2006}.

We have concentrated on numeral phrases; according to Wechsler \& \citet[133-140]{Zlatić2003} similar issues of inflectability depending on the case value arise with uninflectable \textit{nouns} in Serbo-Croat. (There are relatively few of these, but female names such as \textit{Dolôres}, which do not match the appropriate Serbo-Croat inflectional pattern, are uninflectable, as noted in \sectref{sec:key:3.2}.) Ivić (1972: 107n3) pointed out that their use in an oblique case is facilitated by the presence of a modifier or a preposition. See also Bošković (2006, 2009), \citet{Horvath2014} and \citet[103-111]{Adamson2019} on this; there is speaker variation here. External effects have also been suggested for uninflectable \textit{adjectives}; Serbo-Croat also has some instances, many of them orientalisms, as documented in \citet{Nikolić1996}. Bošković claims that these have syntactic effects (2013: 22), notably that an uninflectable adjective must be adjacent to its head noun. His account is discussed in \citet[115-121]{Adamson2019}. There is speaker variation here too.\footnote{ \textrm{\citet[121-131]{Adamson2019} discusses Italian uninflectable adjectives, and concludes that their inability to appear in prenominal position, when unmodified, is not a direct result of their lack of inflection. Hence this is not a case for inclusion here.} } \citet[160-184]{Adamson2019} further analyses German prenominal uninflectable adjectives and argues that here there is a syntactic effect in that they do not allow noun phrase ellipsis.

\subsubsection{\textbf{Definiteness} \textbf{in} \textbf{Bulgarian} }

The relevance of definiteness in Bulgarian was first pointed out by \citet[150-153]{Halpern1995}; see also \citealt{SpencerLuís2012}: 128-130); it is of ongoing interest, as shown by Adamson (2019: 83-103, 2022) and \citet{Georgieva2023}. In Bulgarian, definiteness is inflectional, and a nominal phrase can be marked definite on the first markable item:\footnote{ \textrm{I thank Vladimir Zhobov for his help with the Bulgarian data.} }

Bulgarian (bul, \citealt{SpencerLuís2012}: 128-130)

\ea%26
    \label{ex:key:26}
    \gll \\
         \\
    \glt
    \z % you might need an extra \z if this is the last of several subexamples

          knig-a-ta

  book(\textsc{f)-sg-def.sg.f}

  ‘the book’ 

\ea%27
    \label{ex:key:27}
    \gll \\
         \\
    \glt
    \z % you might need an extra \z if this is the last of several subexamples

          xubav-a-ta    knig-a

  nice-\textsc{sg.f-def.sg.f}  book\textsc{(f)-sg} 

  ‘the nice book’

Example \REF{ex:key:26} shows definiteness marking on the noun since it is the only item in the phrase, while in \REF{ex:key:27} it is on the adjective, as the first markable item. The issue arises where we have an item like \textit{serbez} ‘quarrelsome’, which the researchers cited analyse as an uninflectable adjective. When the phrase is indefinite, all is well, as in \REF{ex:key:28}:

Bulgarian (\citealt{Halpern1995}: 150-153, \citealt{SpencerLuís2012}: 128-130)

\ea%28
    \label{ex:key:28}
    \gll \\
         \\
    \glt
    \z % you might need an extra \z if this is the last of several subexamples

          serbez    žen-a

  quarrelsome  woman\textsc{(f)-sg}

  ‘a quarrelsome woman’

The issue arises when the phrase is definite. If \textit{serbez} ‘quarrelsome’ counts as the first element (as we would expect given \REF{ex:key:27}), ungrammaticality results (the alternatives show the ways that definiteness might, theoretically, be shown on it): 

\ea%29
    \label{ex:key:29}
    \gll \\
         \\
    \glt
    \z % you might need an extra \z if this is the last of several subexamples

          *serbez-ăt/ta    žen-a

  quarrelsome-\textsc{def.sg}   woman\textsc{(f)-sg}

  ‘the quarrelsome woman’

Clearly, then, \textit{serbez} ‘quarrelsome’ does not count as the first element for definiteness. But if it is not counted, the result is also ungrammatical:

\ea%30
    \label{ex:key:30}
    \gll \\
         \\
    \glt
    \z % you might need an extra \z if this is the last of several subexamples

          *serbez  žen-a-ta\footnote{\textrm{Halpern states (1995: 165n22) that this example is possible as an ‘ironic or humorous quotation’ of a previous description.}}

  quarrelsome  woman\textsc{(f)}{}-\textsc{sg-def.sg.f}

  ‘the quarrelsome woman’

Thus neither way works: uninflectable \textit{serbez} ‘quarrelsome’ neither accepts definiteness marking nor allows it on the noun; it neither counts as the first element nor allows the noun to be the first element. Hence we do appear to have an instance where uninflectability has an external effect on the syntax.

\subsection{\textbf{External} \textbf{(syntactic)} \textbf{consequences:} \textbf{summing} \textbf{up}}

I have presented three possible instances and noted others; South Slavonic was well represented. Each case is challenging, and there is speaker variation to be reckoned with. In this section, the key issue is not whether the item inflects, but how similar it is syntactically to regular inflecting items. The Serbo-Croat numeral phrases (\sectref{sec:key:6.2.2}) are significant, since they show defectiveness, in that they fail to fit into syntactic slots where a dative would be required.\footnote{ \textrm{\citet{Testelets2010} offers interesting discussion of case deficiency more generally, based on Russian data.}} 

\section{\textbf{Conclusions}}

We have explored the logical space of uninflectability. First, we calibrated uninflectability across items, determining which items have the property, and how they can be specified. Second, we determined the extent of uninflectability within items. Third, we noted the possible “incoming” conditions on uninflectability. And fourth, we evaluated the “outgoing” requirements imposed by uninflectable items. 

What then can we learn when the dog doesn’t bark, when an item doesn’t inflect? This can challenge our assumptions, and take us forward, in four related areas. First, there is much to be gained for \textbf{morphological} \textbf{typology}. We have seen that in similar morphological systems, uninflectability can vary substantially; recall the way in which Russian has favoured uninflectable items, particularly nouns, while the related Serbo-Croat accepts uninflectable items with reluctance. There are also considerable differences across systems: in more familiar languages we commonly find a minority of items being uninflectable, but we have now seen several languages in which a majority can be uninflectable (whether completely, or in respect of specific features). Given the reasonable assumption that syntactic and morphological classes will align, mismatches between the two are of interest; this is equally true whether the majority of the members of a class are inflectable or uninflectable. And individual items can be very different, both in being fully uninflectable or allowing inflection as an alternative, and in showing various splits in inflectability. In all this, \textbf{diachrony} is a valuable window into how uninflectability is determined; for instance, loss of inflection can reflect a change of part of speech. And we should recall that change is not always from uninflectable items (often the result of borrowing) gaining inflection as they are integrated into the morphological system; we also saw instances of borrowings which gained inflection subsequently losing it. Uninflectability is significant for \textbf{morphological} \textbf{theory}. \citet[143]{Spencer2020} argues that uninflectability itself makes a nonsense of some variants of Distributed Morphology. More generally, it is key to a better understanding of lexical insertion \citep{Spencer2025}. Further, as we have seen, uninflectability raises the question of how inflectional morphology can have external consequences. Thus uninflecting items (and constructions such as Serbo-Croat numeral phrases) may be partially defective, giving unexpected syntactic effects. And finally, uninflectability leads us to question the \textbf{role} \textbf{of} \textbf{inflection}. We saw instances (Russian, Mian, Ingush) where uninflectability appears to make no difference. Similarly, when we examine the complex agreement rules of Archi but find that many agreement targets fail to inflect, this is perplexing. It appears that speakers invest a good deal, in syntactic terms, to ensure the appropriate feature specification for governees and agreement targets, and yet there may be no morphological realization. Why keep a dog when it does not bark? The lack of a bark told Holmes a good deal. Perhaps we need to revisit our assumptions about inflection; its function cannot be the straightforward conveying of grammatical information since it may fail to realize that information, apparently without consequence. 
\begin{verbatim}%%move bib entries to  localbibliography.bib
@misc{Adamson2019,
	author = {Adamson, Luke James},
	note = {} PhD dissertation, University of Pennsylvania. \url{https://repository.upenn.edu/edissertations/3587}{.}},
	title = {\textit{Derivational trapping and the morphosyntax of inflectionlessness},
	year = {2019}
}


@misc{Adamson2022,
	author = {Adamson, Luke},
	note = {\textit{Journal of Slavic Linguistics} 30 (FASL 29 extra issue). 1-13. \url{https://ojs.ung.si/index.php/JSL/article/view/78}.},
	title = {Transparency of inflectionless modifiers for {Bulgarian} definite marker placement},
	year = {2022}
}


@article{Arkadiev2013,
	author = {Arkadiev, Peter M.},
	journal = {\textit{Linguistic Typology}},
	note = {. doi: 10.1515/lingty-2013-0020.},
	pages = {397–437},
	title = {Marking of subjects and objects in {Lithuanian} non-finite clauses: {{A}} typological and diachronic perspective},
	volume = {17},
	year = {2013}
}


@article{Arkadiev2020,
	author = {Arkadiev, Peter M.},
	journal = {\textit{Linguistics}},
	note = {. \href{https://doi.org/10.1515/ling-2020-0045}{{doi: 10.1515/ling-2020-0045}}.},
	number = {2},
	pages = {379–424},
	title = {(Non)finiteness, constructions and participles in {Lithuanian}},
	volume = {58},
	year = {2020}
}


@incollection{Auty1968,
	address = {New York},
	author = {Auty, Robert},
	booktitle = {\textit{Studies in {Slavic} linguistics and poetics in honor of Boris O. {{{U}}}nbegaun}},
	editor = {Robert Magidoff, George Y. Shevelov, J. S. G. Simmons, Kiril Taranovski and John E. Allen},
	pages = {11--14},
	publisher = {New York University Press},
	title = {A note on the indeclinable nouns of Modern {Czech}},
	year = {1968}
}


@book{BaermanBaerman2005,
	address = {Cambridge},
	author = {Baerman, Matthew, Dunstan Brown and Greville G. Corbett},
	number = {109},
	publisher = {Cambridge University Press},
	series = {Cambridge Studies in Linguistics},
	title = {\textit{The syntax-morphology interface: {{A}} study of syncretism}},
	year = {2005}
}


@book{BaermanBaerman2017,
	address = {Cambridge},
	author = {Baerman, Matthew, Dunstan Brown and Greville G. Corbett},
	number = {153},
	publisher = {Cambridge University Press},
	series = {Cambridge Studies in Linguistics},
	title = {\textit{Morphological complexity}},
	year = {2017}
}


Baerman, Matthew \& Greville G. Corbett. 2010. Introduction: defectiveness: typology and diachrony. In: Matthew Baerman, Greville G. Corbett \& Dunstan Brown (eds) \textit{Defective paradigms: missing forms and what they tell us} (Proceedings of the British Academy, 163), 1-18. Oxford: British Academy and Oxford University Press.

@book{Benson1971,
	address = {Philadelphia},
	author = {Benson, Morton with Biljana Šljivić-Šimšić},
	publisher = {University of Pennsylvania Press},
	sortname = {Benson, Morton with Biljana Sljivic-Simsic},
	title = {\textit{{SerboC}roatian-{English} dictionary}},
	year = {1971}
}


@book{BonamiBonami2023,
	address = {\textit{Proceedings of the Society for Computation in Linguistics} 6, Article 29. doi: \href{../../../../../../../../../../Users/lisa/Downloads/10.7275/vyvr-ah62}{{10.7275/vyvr-ah62}}. Available at},
	author = {Bonami, Olivier, Lukáš Kyjánek and Marine Wauquier},
	note = {umass.edu/scil/vol6/iss1/29}.},
	publisher = {\url{https://scholarworks},
	sortname = {Bonami, Olivier, Lukas Kyjanek and Marine Wauquier},
	title = {Assessing the featural organisation of paradigms with distributional methods},
	year = {2023}
}


@incollection{Bond2013,
	address = {Oxford},
	author = {Bond, Oliver},
	booktitle = {\textit{Canonical morphology and syntax}},
	editor = {Dunstan Brown, Marina Chumakina, and Greville G. Corbett},
	pages = {20–47},
	publisher = {Oxford University Press},
	title = {A base for canonical negation},
	year = {2013}
}


@article{Bošković2006,
	author = {Bošković, Željko},
	journal = {\textit{Linguistic Inquiry}},
	number = {3},
	pages = {522–533},
	sortname = {Boskovic, Zeljko},
	title = {Case checking versus case assignment and the case of adverbial {NP}s},
	volume = {37},
	year = {2006}
}


@incollection{Bošković2009,
	address = {Bloomington IN.},
	author = {Bošković, Željko},
	booktitle = {\textit{A linguist’s linguist: {{S}}tudies in {South} {Slavic} linguistics in honor of {E}. {{W}}ayles Browne}},
	editor = {Steven Franks, Vrinda Chidambaram and Brian Joseph},
	pages = {99–122},
	publisher = {Slavica},
	sortname = {Boskovic, Zeljko},
	title = {On \textit{Leo Tolstoy}, its structure, case, left-branch extraction, and prosodic inversion},
	year = {2009}
}


@incollection{Bošković2013,
	address = {Ann Arbor},
	author = {Bošković, Željko},
	booktitle = {\textit{Formal Approaches to {Slavic} Linguistics: {{T}}he Third Indiana \citealt{Meeting2012}}},
	editor = {Steven Franks, Markus Dickinson, George Fowler, Melissa Witcombe and Ksenia Zanon},
	pages = {1–25},
	publisher = {Michigan Slavic Publications},
	sortname = {Boskovic, Zeljko},
	title = {Adjectival escapades},
	volume = {21},
	year = {2013}
}


@incollection{BrownBrown2013,
	address = {Oxford},
	author = {Brown, Dunstan and Marina Chumakina},
	booktitle = {\textit{Canonical morphology and syntax}},
	editor = {Dunstan Brown, Marina Chumakina and Greville G. Corbett},
	pages = {1--19},
	publisher = {Oxford University Press},
	title = {What there might be and what there is: {{A}}n introduction to Canonical Typology},
	year = {2013}
}


@book{BrownEtAl2007,
	address = {The alignment of form and function},
	author = {Brown, Dunstan, Carole Tiberius and and Corbett},
	note = {\textit{International Journal of Corpus Linguistics} 12.511\nobreakdash-534.},
	publisher = {Corpus-based evidence from Russian},
	year = {2007}
}


@article{BrownBrown2008,
	author = {Brown, Lea and Matthew S. Dryer},
	journal = {\textit{Language}},
	pages = {528–565},
	title = {The verbs for ‘and’ in {Walman}, a {Torricelli} language of {Papua New Guinea}},
	volume = {84},
	year = {2008}
}


Browne, Wayles. 1993a. Serbo-Croat. In: Bernard Comrie \& Greville G. Corbett (eds) \textit{The Slavonic languages}, 306-387. London: Routledge.

@article{Browne1993b,
	author = {Browne, Wayles},
	journal = {\textit{Zbornik Matice srpske za filologiju i lingvistiku}},
	number = {2},
	pages = {173--179},
	title = {Variation in {Serbo-Croatian} noun phrase constituents: {{D}}eclinable and indeclinable items},
	volume = {36},
	year = {1993}
}


@book{Browne2015,
	address = {\textit{Balkanistica} 28 (=\textit{Od Čikago i nazad},
	author = {Browne, Wayles},
	note = {Friedman on the occasion of his retirement}, edited by Donald L. Dyer, Christina E. Kramer \& Brian D. Joseph), 65–78.},
	publisher = {Papers to Honor Victor A},
	title = {The noun strikes back},
	year = {2015}
}


@incollection{Browne2022,
	address = {Berlin},
	author = {Browne, Wayles},
	booktitle = {\textit{Lifetime linguistic inspirations: {{T}}o Igor Mel’čuk from colleagues and friends for his 90th birthday} (\textit{{Wiener} Slawistischer Almanach}, Linguistische {Reihe}, Sonderband 101)},
	editor = {Leonid Iomdin, Jasmina Milićević and Alain Polguère},
	pages = {129--136},
	publisher = {Peter Lang},
	title = {Premodifiers on {Bosnian}-{Croatian}-{Serbian} “numeral + noun” phrases},
	year = {2022}
}


@article{Bugarski2012,
	author = {Bugarski, Ranko},
	journal = {\textit{Journal of Multilingual and Multicultural Development}},
	note = {. \href{https://doi.org/doi:%/01434632.2012.663376}{doi: 10.1080/01434632.2012.663376}.},
	pages = {219--235},
	title = {Language, identity and borders in the former {Serbo-Croatian} area},
	urldate = {2010.1080},
	volume = {33},
	year = {2012}
}


@incollection{Bugarski2019,
	address = {Budapest},
	author = {Bugarski, Ranko},
	booktitle = {\textit{Language variation: {{A}} factor of increasing complexity and a challenge for language policy within {Europe}: {{C}}ontributions to the {EFNIL} \citealt{Conference2018} in {Amsterdam}}},
	editor = {Tanneke Schoonheim and Johan Van Hoorde},
	note = {\url{},
	pages = {105--114},
	publisher = {Research Institute for Linguistics, Hungarian Academy of Sciences},
	title = {Past and current developments involving pluricentric {Serbo-Croatian} and its official heirs},
	url = {http://www.efnil.org/documents/conference-publications/amsterdam-2018}},
	year = {2019}
}


@incollection{Bye2015,
	address = {London},
	author = {Bye, Patrik},
	booktitle = {\textit{Understanding allomorphy: {{P}}erspectives from optimality theory}},
	editor = {Eulàlia Bonet, Maria-Rosa Lloret and Joan Mascaró},
	pages = {107--176},
	publisher = {Equinox},
	title = {The nature of allomorphy and exceptionality: {{{E}}}vidence from {Burushaski} plurals},
	year = {2015}
}


@article{Cetnarowska2021,
	author = {Cetnarowska, Bożena},
	journal = {\textit{Word Structure}},
	note = {. doi: 10.3366/word.2021.0180.},
	number = {1},
	pages = {59–96},
	sortname = {Cetnarowska, Bozena},
	title = {Variability in inflectional forms of attributive-appositive composite lexical units in {Polish}},
	volume = {14},
	year = {2021}
}


@article{ChuangChuang2022,
	author = {Chuang, Yu-Ying, Dunstan Brown, R. Harald Baayen and Roger Evans},
	journal = {\textit{The Mental Lexicon}},
	note = {. doi: \href{http://dx.doi.org/10.1075/ml.22013.chu}{10.1075/ml.22013.chu}.},
	number = {3},
	pages = {395–421},
	title = {Paradigm gaps are associated with weird “distributional semantics” properties: {{R}}ussian defective nouns and their case and number paradigms},
	volume = {17},
	year = {2022}
}


@book{Chumakina2023,
	address = {\textit{Mediterranean Morphology Meetings: 13: Comparing typologies}, 37-46. At},
	author = {Chumakina, Marina},
	note = {library.upatras.gr/mmm/article/view/4404/4325.},
	publisher = {https://pasithee},
	title = {Locative forms in {Nakh-Daghestanian} as an example of a transcategorial paradigm},
	year = {2023}
}


Chumakina, Marina \& Oliver Bond. 2016. Competing controllers and agreement potential. In Oliver Bond, Greville G. Corbett, Marina Chumakina \& Dunstan Brown (eds) \textit{Archi: Complexities of agreement in cross-theoretical perspective,} 77-117. Oxford: Oxford University Press.

@incollection{ChumakinaChumakina2015,
	address = {Oxford},
	author = {Chumakina, Marina and Greville G. Corbett},
	booktitle = {\textit{Understanding and measuring morphological complexity}},
	editor = {Matthew Baerman, Dunstan Brown and Greville G. Corbett},
	pages = {93--116},
	publisher = {Oxford University Press},
	title = {Gender-number marking in {Archi}: {{S}}mall is complex},
	year = {2015}
}


@book{ChuprinkoChuprinko2023,
	address = {\textit{Acta Linguistica Academica}. doi},
	author = {Chuprinko, Kirill, Varvara Magomedova and Natalia Slioussar},
	note = {.2023.00641.},
	publisher = {10},
	title = {Gender variation in indeclinable inanimate nouns and gender markedness in modern {Russian}},
	urldate = {1556/2062},
	year = {2023}
}


@book{ComrieComrie1993,
	address = {London},
	booktitle = {\textit{The {Slavonic} Languages}},
	editor = {Comrie, Bernard and Greville G. Corbett},
	publisher = {Routledge},
	title = {\textit{The {Slavonic} Languages}},
	year = {1993}
}


@book{ComrieComrie1996,
	address = {\textit{The} \textit{Russian language in the Twentieth Century.} Oxford},
	author = {Comrie, Bernard, Gerald Stone and Maria Polinsky},
	publisher = {Clarendon Press},
	year = {1996}
}


Copot, Maria, Ninoh Agostinho Da Silva, Ahmed Beji, Arno Watiez \& Olivier Bonami. 2025. Emerging uninflectedness in French clipped verbs. In Enrique Palancar \& Sebastian Fedden (eds) \textit{Uninflectedness}. Berlin: Language Science Press.

@article{Corbett1987,
	author = {Corbett, Greville G.},
	journal = {\textit{Language}},
	pages = {299--345},
	title = {The morphology/syntax interface: {{{E}}}vidence from possessive adjectives in {Slavonic}},
	volume = {63},
	year = {1987}
}


@book{Corbett1991,
	address = {Cambridge},
	author = {Corbett, Greville G.},
	publisher = {Cambridge University Press},
	title = {\textit{Gender}},
	year = {1991}
}


@book{Corbett2006,
	address = {Cambridge},
	author = {Corbett, Greville G.},
	publisher = {Cambridge University Press},
	title = {\textit{Agreement}},
	year = {2006}
}


Corbett, Greville G. 2009. Morphology-free syntax: two potential counter-examples from Serbo-Croat. In: Steven Franks, Vrinda Chidambaram \& Brian Joseph (eds) \textit{A Linguist's Linguist: Studies in South Slavic linguistics in honor of E. Wayles Browne}, 149\nobreakdash-166. Bloomington, Indiana: Slavica.

@incollection{Corbett2011,
	address = {Berlin},
	author = {Corbett, Greville G.},
	booktitle = {\textit{{E}xpecting the unexpected: {{{E}}}xceptions in grammar} (Trends in Linguistics. {{S}}tudies and Monographs 216)},
	editor = {Horst J. Simon and Heike Wiese},
	pages = {107--126},
	publisher = {De Gruyter Mouton},
	title = {Higher order exceptionality in inflectional morphology},
	year = {2011}
}


@book{Corbett2012,
	address = {Cambridge},
	author = {Corbett, Greville G.},
	publisher = {Cambridge University Press},
	title = {\textit{Features}},
	year = {2012}
}


@incollection{Corbett2013,
	address = {Oxford},
	author = {Corbett, Greville G.},
	booktitle = {\textit{Canonical morphology and syntax}},
	editor = {Dunstan Brown, Marina Chumakina and Greville G. Corbett},
	pages = {48--65},
	publisher = {Oxford University Press},
	title = {Canonical morphosyntactic features},
	year = {2013}
}


@article{Corbett2015,
	author = {Corbett, Greville G.},
	journal = {\textit{Language}},
	pages = {145--193},
	title = {Morphosyntactic complexity: {{A}} typology of lexical splits},
	volume = {91},
	year = {2015}
}


@article{Corbett2019,
	author = {Corbett, Greville G.},
	journal = {\textit{Morphology}},
	note = {. doi: 10.1007/s11525-018-9336-0.},
	pages = {51–108},
	title = {Pluralia tantum nouns and the theory of features: {{A}} typology of nouns with non-canonical number properties},
	volume = {29},
	year = {2019}
}


@book{Corbett2023,
	address = {\textit{Glossa: a journal of general linguistics} 8(1). doi},
	author = {Corbett, Greville G.},
	note = {org/10.16995/glossa.9164}{{10.16995/glossa.9164}}.},
	publisher = {\href{https://doi},
	title = {The Agreement Hierarchy and (generalized) semantic agreement},
	year = {2023}
}


@incollection{CorbettCorbett2018,
	address = {London},
	author = {Corbett, Greville G. and Wayles Browne},
	booktitle = {\textit{The World’s major languages,} 3\textsuperscript{rd} edition},
	editor = {Bernard Comrie},
	pages = {339--356},
	publisher = {Routledge},
	title = {Serbo-{Croat}: {{B}}osnian, {Croat}ian, Montenegrin, {Serbian}},
	year = {2018}
}


@article{CorbettCorbett2016,
	author = {Corbett, Greville G. and Sebastian Fedden},
	journal = {\textit{Journal of Linguistics}},
	pages = {495–531},
	title = {Canonical gender},
	volume = {52},
	year = {2016}
}


@article{CorbettCorbett2017,
	author = {Corbett, Greville G., Sebastian Fedden and Raphael Finkel},
	journal = {\textit{Linguistic Typology}},
	pages = {209–260},
	title = {Single versus concurrent feature systems: {{N}}ominal classification in {Mian}},
	volume = {21},
	year = {2017}
}


@incollection{D’AchilleD’Achille2003,
	address = {Rome},
	author = {D’Achille, Paolo and Anna M. Thornton},
	booktitle = {\textit{{{Italia}} linguistica anno Mille – {{Italia}} linguistica anno Duemila.} \textit{{Atti} del {XXXIV} Congresso della Società di Linguistica {{Italia}}na}},
	editor = {Nicoletta Maraschio and Teresa Poggi Salani},
	pages = {211–230},
	publisher = {Bulzoni},
	sortname = {D'Achille, Paolo and Anna M. Thornton},
	title = {{La} flessione del nome dall’italiano antico all’italiano contemporaneo},
	year = {2003}
}


@misc{Doleschal2000,
	author = {Doleschal, Ursula},
	note = {\textit{Das Phänomen der Unflektierbarkeit in den slawischen Sprachen.} Habilitationsschrift. \citealt{Vienna2000} [improved version 2001].},
	title = {-2001},
	year = {2000}
}


@incollection{Doleschal2002a,
	address = {Prague},
	author = {Doleschal, Ursula},
	booktitle = {\textit{Setkání s češtinou}},
	editor = {Alena Krausová, Markéta Slezáková and Zdeňka Svobodová},
	pages = {8--12},
	publisher = {Ústav pro jazyk český AV ČR},
	title = {Aspekty ohebnosti a neohebnosti substantiv v češtině},
	year = {2002}
}


@incollection{Doleschal2002b,
	address = {Munich},
	author = {Doleschal, Ursula},
	booktitle = {\textit{Linguistische {Beiträge} zur Slavistik aus {Deutschland} und Österreich. (Jungslavistinnentreffen \citealt{Wittenberg2000})}},
	editor = {Thomas Daiber},
	pages = {55--66},
	publisher = {Otto Sagner},
	title = {“Super sexy Mici” Oder das {Phänomen} Der Unflektierbarkeit in den slawischen {Sprachen}},
	year = {2002}
}


@incollection{Doleschal2008,
	address = {Munich},
	author = {Doleschal, Ursula},
	booktitle = {\textit{Linguistische {Beiträge} zur Slavistik. {{X}}{IV}. {{J}}ungslavistInnen-Treffen in Stuttgart 15.-18. \citealt{September2005}}. (Specimina Philologiae Slavicae; )},
	editor = {Ljudmila Geist and Grit Mehlhorn},
	pages = {51--64},
	publisher = {Otto Sagner},
	title = {Über die {Entstehung} des Deklinationstyps \textit{kuli, kuliho}},
	volume = {150},
	year = {2008}
}


Doleschal, Ursula. 2025. The conditions of uninflectedness in nouns in the Slavic languages. In Enrique Palancar \& Sebastian Fedden (eds) \textit{Uninflectedness}. Berlin: Language Science Press.

@article{Dunn1988,
	author = {Dunn, John A.},
	journal = {Is there a tendency towards analyticity in Russian? \textit{The Slavonic and East European Review}},
	number = {2},
	pages = {169--183},
	volume = {66},
	year = {1988}
}


@incollection{Dunn1989,
	address = {Nottingham},
	author = {Dunn, John A.},
	booktitle = {\textit{Words and images: {{{E}}}ssays in honour of Professor (emeritus) Dennis Ward}},
	editor = {Michael Falchikov, Christopher Pike and Robert Russell},
	pages = {81--95},
	publisher = {Astra Press},
	title = {Notes on indeclinable nouns in {Russian}},
	year = {1989}
}


@article{Ďurovič1985,
	author = {Ďurovič, Ľubomír},
	journal = {\textit{International Journal of Slavic Linguistics and Poetics}},
	pages = {0538--8228},
	sortname = {Durovic, Lubomir},
	title = {The numerals in {Serbo-Croatian}},
	volume = {31-32.123-130},
	year = {1985}
}


@book{Emeneau1955,
	address = {12). Berkeley},
	author = {Emeneau, Murray B.},
	publisher = {University of California Press},
	title = {\textit{Kolami: {{A}} {Dravidian} language} (University of {California} Publications in Linguistics vol},
	year = {1955}
}


Esher, Louise. 2025. Loss of inflection in the diachrony of French nouns. In Enrique Palancar \& Sebastian Fedden (eds) \textit{Uninflectedness}. Berlin: Language Science Press.

@incollection{Fedden2019,
	address = {Edinburgh},
	author = {Fedden, O. Sebastian},
	booktitle = {\textit{Morphological perspectives: {{P}}apers in honour of Greville G. {{C}}orbett}},
	editor = {Matthew Baerman, Oliver Bond and Andrew Hippisley},
	pages = {303--326},
	publisher = {Edinburgh University Press},
	title = {{To} agree or not to agree? {{A}} typology of sporadic agreement},
	year = {2019}
}


@book{Fedden2022,
	address = {\textit{Word Structure} 15.3 (special issue \textit{The many facets of agreement}, ed. by Tania Paciaroni, Alice Idone and Michele Loporcaro), 283-304. doi},
	author = {Fedden, Sebastian},
	note = {3366/word.2022.0211.},
	publisher = {10},
	title = {Agreement and argument realization in {Mian} discourse},
	year = {2022}
}


Fedden, Sebastian, Matías Guzmán Naranjo \& Greville G. Corbett. Typology meets statistical modeling: the German gender system. To appear in \textit{Language}.

@book{FeddenFedden2022,
	address = {Reduced agreement},
	author = {Fedden, Sebastian and Tania Paciaroni},
	note = {Paper at \citealt{SLE2022}, 55\textsuperscript{th} Annual Meeting, Bucharest. .},
	publisher = {A typological approach},
	urldate = {26.08.2022},
	year = {2022}
}


@article{Forker2018,
	author = {Forker, Diana},
	journal = {\textit{Linguistics}},
	pages = {865--894},
	title = {Gender agreement is different},
	volume = {56},
	year = {2018}
}


Franks, Steven. Forthcoming. Numerals and quantity expressions. In Danko Šipka \& Wayles Browne (eds) \textit{The Cambridge handbook of Slavic linguistics}. Cambridge: Cambridge University Press.

Friedewald, Katja. 2025. French \textit{voilà}: An uninflectable form arising from an inflecting verb. In Enrique Palancar \& Sebastian Fedden (eds) \textit{Uninflectedness}. Berlin: Language Science Press.

Gaszewski, Jerzy. 2025. Uninflectedness as a rule in Polish, an inflected language. In Enrique Palancar \& Sebastian Fedden (eds) \textit{Uninflectedness}. Berlin: Language Science Press.

@incollection{Georgieva2023,
	address = {Berlin: Language Science Press. doi},
	author = {Georgieva, Ekaterina},
	booktitle = {\textit{Advances in formal {Slavic} linguistics 2021}},
	editor = {Petr Biskup, Marcel Börner, Olav Mueller-Reichau and Iuliia Shcherbina},
	note = {5281/zenodo.10123637.},
	pages = {179–209},
	publisher = {10},
	title = {Inflectionless adjectives in {Bulgarian} as a case of nominal predication},
	year = {2023}
}


@article{Giger2016,
	author = {Giger, Markus},
	journal = {\textit{Jazykovedný časopis}},
	note = {. doi: \url{https://doi.org/10.1515/jazcas-2017-0013}.},
	number = {3},
	pages = {283--295},
	title = {Kongruenzbrüche in slovakischen possessiven Resultativa (Evidenz aus dem {SlovakischenN}ationalkorpus)},
	volume = {67},
	year = {2016}
}


@article{Grubišić1994,
	author = {Grubišić, Vinko},
	journal = {\textit{Jezik}},
	note = {. \url{https://hrcak.srce.hr/252611}.},
	number = {3},
	pages = {78–82},
	sortname = {Grubisic, Vinko},
	title = {Mijenjaju li se brojevi \textit{dva}, \textit{tri} i \textit{četiri} kad su s prijedlozima},
	volume = {42},
	year = {1994}
}


@book{Halpern1995,
	address = {\textit{On the placement and morphology of clitics.} Stanford},
	author = {Halpern, Aaron},
	publisher = {CSLI Publications},
	year = {1995}
}


@article{Hamm1956,
	author = {Hamm, Josip},
	journal = {\textit{Jezik}},
	note = {. \url{https://hrcak.srce.hr/clanak/102015}.},
	number = {1},
	pages = {9--14},
	title = {Promjena brojeva 2, 3, i 4},
	volume = {5},
	year = {1956}
}


Holisky, Dee Ann \& Rusudan Gagua. 1994. Tsova–Tush (Batsbi). In Rieks Smeets (ed.) \textit{Indigenous languages of the Caucasus IV: North East Caucasian languages II: presenting the three Nakh Languages and six minor Lezgian languages,} 147–212. Delmar, NY: Caravan Books.

@incollection{Horvath2014,
	address = {Amsterdam},
	author = {Horvath, Julia},
	booktitle = {\textit{Advances in the syntax of {DP}s: {{S}}tructure, agreement, and case}},
	editor = {Anna Bondaruk, Gréte Dalmi and Alexander Grosu},
	pages = {117–128},
	publisher = {John Benjamins},
	title = {A note on oblique case: {{{E}}}vidence from {Serbian}/{Croatian}},
	year = {2014}
}


@book{Ivić1972,
	address = {\textit{Russkoe i slavjanskoe jazykoznanie: K 70-letiju člena-korrespondenta Akademii nauk SSSR R. I. Avanesova,} 106-21. Moscow},
	author = {Ivić, Pavle},
	publisher = {Nauka},
	sortname = {Ivic, Pavle},
	title = {Sistema padežnyx okončanij suščestvitel′nyx v serboxorvatskom literaturnom jazyke},
	year = {1972}
}


@book{Kibrik1977,
	address = {\textit{Opyt strukturnogo opisanija arčinskogo jazyka: III: Dinamičeskaja grammatika.} Moscow},
	author = {Kibrik, Aleksandr E.},
	note = {Moskovskogo universiteta.},
	publisher = {Izd},
	year = {1977}
}


Kibrik, Aleksandr E. 1994. Archi. In Rieks Smeets (ed.) \textit{Indigenous languages of the Caucasus IV: North East Caucasian languages II: presenting the three Nakh Languages and six minor Lezgian languages,} 297-365. Delmar, NY: Caravan Books.

@book{KibrikKibrik1990,
	address = {Moscow},
	author = {Kibrik, \textcyrillic{А}leksandr. \textcyrillic{Е}. and Sandro V. Kodzasov},
	publisher = {MGU},
	sortname = {Kibrik, \textcyrillic{А}leksandr. \textcyrillic{Е}. and Sandro V. Kodzasov},
	title = {\textit{Sopostavitel′noe izučenie dagestanskix jazykov: {{I}}mja: {{F}}onetika}},
	year = {1990}
}


Köhlich, Viktor. 2025. The uninflecting word class \textit{rentaishi} in Modern Japanes. In Enrique Palancar \& Sebastian Fedden (eds) \textit{Uninflectedness}. Berlin: Language Science Press.

@incollection{Koptev2008,
	address = {Helsinki},
	author = {Koptev, Mikhail V.},
	booktitle = {\textit{Instrumentarij rusistiki: {{K}}orpusnye podkhody} (Slavica Helsingiensia 34)},
	editor = {L. A. Birjulin, Mikhail Kopotev, Arto Mustajoki and Ekaterina J. Protasova},
	pages = {136--151},
	publisher = {University of Helsinki},
	title = {K postroeniju častotnoj grammatiki russkogo jazyka: {{P}}adežnaja Sistema po korpusnym dannym},
	year = {2008}
}


@article{Kovačević1951,
	author = {Kovačević, Srbislava},
	journal = {\textit{Naš jezik}, n.s.},
	number = {7-10},
	pages = {246--255},
	sortname = {Kovacevic, Srbislava},
	title = {O imenici \textit{doba} u našem jeziku},
	volume = {2},
	year = {1951}
}


@article{Kraljević1976,
	author = {Kraljević, Gojko},
	journal = {\textit{Naš jezik}, n.s.},
	note = {. \url{https://dais.sanu.ac.rs/123456789/8076}.},
	number = {1-2},
	pages = {43--52},
	sortname = {Kraljevic, Gojko},
	title = {Deklinabilnost i indeklinabilnost brojeva \textit{dva}, \textit{tri}, \textit{četiri}},
	volume = {22},
	year = {1976}
}


@article{Krzyżanowski2020,
	author = {Krzyżanowski, Piotr},
	journal = {\textit{Nieodmienność rzeczowników w systemie norm współczesnej.} \href{https://www.ceeol.com/search/journal-detail?id=1043}{{\textit{Język Polski}}}},
	note = {. \href{https://doi.org/10.31286/JP.100.2.4}{{doi: 10.31286/JP.100.2.4}}.},
	number = {2},
	pages = {49--59},
	sortname = {Krzyzanowski, Piotr},
	volume = {100},
	year = {2020}
}


@article{Kuzmić2007,
	author = {Kuzmić, Boris},
	journal = {do 18. stoljeća. \textit{Croatica et Slavica Iadertina}},
	note = {. doi: 10.15291/csi.357.},
	pages = {9--29},
	sortname = {Kuzmic, Boris},
	title = {Deklinacija brojeva \textit{dva}, \textit{oba}, \textit{tri} i \textit{četiri} u čakavskim pravnim tekstovima od 13},
	volume = {3},
	year = {2007}
}


@article{KuzmićKuzmić2007,
	author = {Kuzmić, Boris and Martina Kuzmić},
	journal = {do 18. stoljeća. \textit{Rasprave Instituta za hrvatski jezik i jezikoslovlje}},
	pages = {263–288},
	sortname = {Kuzmic, Boris and Martina Kuzmic},
	title = {Deklinacija brojeva \textit{dva}, \textit{oba}, \textit{tri} i \textit{četiri} u kajkavskim pravnim tekstovima od 16},
	volume = {33},
	year = {2007}
}


Kyuseva, Maria, Alexander Krasovitsky, Matthew Baerman \& Greville G. Corbett. 2025. Diachronic paths to uninflectedness in South Slavonic. In Enrique Palancar \& Sebastian Fedden (eds) \textit{Uninflectedness}. Berlin: Language Science Press.

@incollection{Loporcaro2011,
	address = {Oxford: Oxford University Press. \href{https://doi.org/10.1093/acprof:oso/9780199589982.003.0016}{doi: 10.1093/acprof},
	author = {Loporcaro, Michele},
	booktitle = {\textit{Morphological autonomy: {{P}}erspectives from {Romance} inflectional morphology}},
	editor = {Martin Maiden, John Charles Smith, Maria Goldbach and Marc-Olivier Hinzelin},
	note = {003.0016}.},
	pages = {327–357},
	publisher = {oso/9780199589982},
	title = {Syncretism and neutralization in the marking of {Romance} object agreement},
	year = {2011}
}


Loporcaro, Michele. 2025. Uninflectedness as a factor in agreement loss. In Enrique Palancar \& Sebastian Fedden (eds) \textit{Uninflectedness}. Berlin: Language Science Press.

@book{Lowe2017,
	address = {\textit{Transactions of the Philological Society} \href{https://onlinelibrary.wiley.com/toc/1467968x/2017/115/2}{\textstyleval{115}{(}\textstyleval{2}}).263-294. \href{https://doi.org/10.1111/1467-968X.12102}{{doi},
	author = {Lowe, John J.},
	note = {-968X.12102}}.},
	publisher = {10},
	title = {The {Sanskrit} (pseudo)periphrastic future},
	urldate = {1111/1467},
	year = {2017}
}


@book{Maldjieva1995,
	address = {Warsaw},
	author = {Maldjieva, Vyara},
	publisher = {Energeia},
	title = {\textit{Non-inflected parts of speech in the {Slavonic} languages}},
	year = {1995}
}


@article{Marković1951,
	author = {Marković, Svetozar},
	journal = {\textit{Naš jezik}, n.s.},
	number = {5-6},
	pages = {155--161},
	sortname = {Markovic, Svetozar},
	title = {O promenljivosti broja \textit{dva}},
	volume = {2},
	year = {1951}
}


Matushansky, Ora. 2022. On the semantics of inanimate gender. In Linnaea Stockall, Luisa Martí, David Adger, Isabelle Roy \& Sarah Ouwayda (eds) \textit{For Hagit: A celebration. חגיגת חגית/xagigat xagit}, \textit{QMUL Occasional Papers in Linguistics} 47. London: QMUL.

@incollection{Matushansky2023,
	address = {Berlin: Language Science Press. doi},
	author = {Matushansky, Ora},
	booktitle = {\textit{Advances in formal {Slavic} linguistics 2021}},
	editor = {Petr Biskup, Marcel Börner, Olav Mueller-Reichau and Iuliia Shcherbina},
	note = {5281/zenodo.10123643.},
	pages = {233–263},
	publisher = {10},
	title = {Phi-congruence and case-agreement in close apposition in {Russian}},
	year = {2023}
}


@incollection{Mayo1993,
	address = {London},
	author = {Mayo, Peter},
	booktitle = {\textit{The {Slavonic} languages}},
	editor = {Bernard Comrie and Greville G. Corbett},
	pages = {887--946},
	publisher = {Routledge},
	title = {Belorussian},
	year = {1993}
}


@article{Mirtov1953,
	author = {Mirtov, A. V.},
	journal = {\textit{Voprosy jazykoznanija} no},
	pages = {99--108},
	title = {Iz nabljudenij nad russkim jazykom v èpoxu Velikoj Otečestvennoj vojny},
	volume = {4},
	year = {1953}
}


@incollection{Mučnik1964,
	address = {Moskva},
	author = {Mučnik, I. P.},
	booktitle = {\textit{Razvitie grammatiki i leksiki sovremennogo russkogo jazyka}},
	editor = {I. P. Mučnik and M. V. Panov},
	pages = {148--180},
	publisher = {Nauka},
	sortname = {Mucnik, I. P.},
	title = {Neizmenjaemye suščestvitel′nye, ix mesto v sisteme sklonenija i tendencii razvitija v sovremennom russkom literaturnom jazyke},
	year = {1964}
}


@misc{Murphy2000,
	author = {Murphy, Dianna L.},
	note = {Doctoral dissertation, Ohio State University. \url{http://rave.ohiolink.edu/etdc/view?acc_num=osu25666262}.},
	title = {The gender of inanimate indeclinable common nouns in Modern {Russian}},
	urldate = {14881948},
	year = {2000}
}


@article{Naylor1982,
	author = {Naylor, Kenneth},
	journal = {\textit{International Journal of Slavic Linguistics and Poetics}},
	pages = {291--295},
	title = {Phonology affecting morphology: {{T}}he case of Serbocroatian indeclinables},
	volume = {25-26},
	year = {1982}
}


@book{Nichols2011,
	address = {\textit{Ingush grammar.} Berkeley},
	author = {Nichols, Johanna},
	note = {\url{https://escholarship.org/uc/item/3nn7z6w5}.},
	publisher = {University of California Press},
	year = {2011}
}


@article{Nichols2018,
	author = {Nichols, Johanna},
	journal = {\textit{Linguistics}},
	pages = {845--863},
	title = {Agreement with overt and null arguments in {Ingush}},
	volume = {56},
	year = {2018}
}


@article{NikolaevNikolaev2022,
	author = {Nikolaev, Alexandre and Neil \citealt{Bermel},
	journal = {Explaining uncertainty and defectivity of inflectional paradigms. \textit{Cognitive Linguistics}},
	note = {. doi: 10.1515/cog-2021-0041.},
	number = {3},
	pages = {585--621},
	sortname = {Nikolaev, Alexandre and Neil \citealt{Bermel},
	title = {}},
	volume = {33},
	year = {2022}
}


@incollection{Nikolaeva2013,
	address = {Oxford},
	author = {Nikolaeva, \citealt{Irina},
	booktitle = {\textit{Canonical morphology and syntax}},
	editor = {Dunstan Brown, Marina Chumakina and Greville G. Corbett},
	pages = {99–122},
	publisher = {Oxford University Press},
	sortname = {Nikolaeva, \citealt{Irina},
	title = {}. {{{U}}}npacking finiteness},
	year = {2013}
}


@article{Nikolić1976,
	author = {Nikolić, Miroslav},
	journal = {\textit{Naš jezik}, n.s.},
	note = {. \url{https://dais.sanu.ac.rs/123456789/8078}.},
	number = {1-2},
	pages = {52--56},
	sortname = {Nikolic, Miroslav},
	title = {Promenljivost brojeva 2-4 u štokavskim govorima},
	volume = {22},
	year = {1976}
}


@misc{Nikolić1995,
	author = {Nikolić, Miroslav},
	note = {Nedeklinabilne imenice u srpskom jeziku. \textit{Naš jezik}, n.s., 30(/1\nobreakdash-5).15-43. \url{https://dais.sanu.ac.rs/123456789/7819}.},
	title = {-96},
	year = {1995}
}


@misc{Nikolić1996,
	author = {Nikolić, Miroslav},
	note = {Naš jezik}, n.s., 31(1\nobreakdash-5).35\nobreakdash-54. \url{https://dais.sanu.ac.rs/123456}.},
	title = {Nepromenljivi pridevi u srpskom jeziku\textit{},
	urldate = {789/1361},
	year = {1996}
}


@article{Nübling2017,
	author = {Nübling, Damaris},
	journal = {\textit{Folia Linguistica}},
	note = {. doi: 10.1515/flin-2017-0012.},
	number = {2},
	pages = {341–367},
	sortname = {Nubling, Damaris},
	title = {The growing distance between proper names and common nouns in {German}: {{O}}n the way to onymic schema constancy},
	volume = {51},
	year = {2017}
}


Palancar, Enrique. 2025. Uninflectedness in Amuzgan verbal inflection. In Enrique Palancar \& Sebastian Fedden (eds) \textit{Uninflectedness}. Berlin: Language Science Press.

@book{Panov1968,
	address = {\textit{Morfologija i sintaksis sovremennogo russkogo literaturnogo jazyka (Russkij jazyk i sovetskoe obščestvo: III).} Moscow},
	editor = {Panov, M. V.},
	publisher = {Nauka},
	year = {1968}
}


@book{Popović1979,
	address = {\textit{Jugoslovenski seminar za strane slaviste} 30.3\nobreakdash-24. Belgrade: Filološki fakultet u Beogradu},
	author = {Popović, Ljubomir},
	publisher = {Međunarodni slavistički centar},
	sortname = {Popovic, Ljubomir},
	title = {Upotreba kardinalnih brojeva u srpskohrvatskom jeziku},
	year = {1979}
}


@article{Popović1981,
	author = {Popović, Ljubomir},
	journal = {Nesuglašenost forme i funkcije nepromenljivih brojeva u srpskohrvatskom jeziku: problem dativa. \textit{Makedonski jazik}},
	note = {. ISSN: .},
	pages = {603--609},
	sortname = {Popovic, Ljubomir},
	title = {-82},
	urldate = {0025-1089},
	volume = {32-33},
	year = {1981}
}


@incollection{Popović1982,
	address = {Warsaw},
	author = {Popović, Ljubomir},
	booktitle = {\textit{Studia gramatyczne V}},
	editor = {Stanisław Urbańczyk},
	pages = {149--152},
	publisher = {Prace Instytutu jezyka Polskiego},
	sortname = {Popovic, Ljubomir},
	title = {Reperkusije neoznačenosti dativa nepromenljivih brojeva na gramatički sistem i na strukturu teksta u srpskohrvatskom jeziku},
	year = {1982}
}


@article{Popović1967,
	author = {Popović, Milenko},
	journal = {\textit{Jezik}},
	number = {5},
	pages = {144--148},
	sortname = {Popovic, Milenko},
	title = {O brojnim konstrukcijama kao blokovima koji se sklanjaju},
	volume = {14},
	year = {1967}
}


@article{Radović-Tešić1976,
	author = {Radović-Tešić, Milica},
	journal = {\textit{Naš jezik} n.s},
	note = {. \url{https://dais.sanu.ac.rs/123456789/3792}.},
	number = {1-2},
	pages = {56--63},
	sortname = {Radovic-Tesic, Milica},
	title = {Neki podaci i stavovi o promenljivosti brojeva},
	volume = {22},
	year = {1976}
}


Renner, Vincent \& Adam Renwick. 2025. Inflection dropping in the English-origin verbs of present-day French: A Twitter-wide exploration. In Enrique Palancar \& Sebastian Fedden (eds) \textit{Uninflectedness}. Berlin: Language Science Press.

Rhodes, \citealt{Richard1987}. Paradigms large and small. In Jon Aske, Natasha Beery, Laura Michaelis \& Hana Filip (eds) \textit{Proceedings of the Thirteenth Annual Meeting of the Berkeley Linguistics Society: General session and parasession on grammar and cognition,} 223-34. Berkeley, California: Berkeley Linguistics Society, University of California.

@article{Rogić1954,
	author = {Rogić, Pavle},
	journal = {\textit{Jezik}},
	number = {5},
	pages = {138--141},
	sortname = {Rogic, Pavle},
	title = {Deklinacija brojeva dva, oba (obadva), tri, četiri},
	volume = {3},
	year = {1954}
}


@book{Round2013,
	address = {Oxford},
	author = {Round, Erich R.},
	publisher = {Oxford University Press},
	title = {\textit{Kayardild morphology and syntax}},
	year = {2013}
}


@book{Round2023,
	address = {\textit{Linguistic Typology} aop. doi},
	author = {Round, Erich R.},
	note = {doi.org/10.1515/lingty-2022-0032}{10.1515/lingty-2022-0032}.},
	publisher = {\href{http://dx},
	title = {Canonical phonology and criterial conflicts: {{R}}elating and resolving four dilemmas of phonological typology},
	year = {2023}
}


@article{RoundRound2020,
	author = {Round, Erich R. and Greville G. Corbett},
	journal = {\textit{Linguistic Typology}},
	pages = {489--525},
	title = {Comparability and measurement in typological science: {{T}}he bright future for linguistics},
	volume = {24},
	year = {2020}
}


@article{SchlechtwegSchlechtweg2023,
	author = {Schlechtweg, Marcel and G. Corbett},
	journal = {Is morphosyntactic agreement reflected in acoustic detail? \textit{English Language and Linguistics}},
	note = {. doi: \href{https://doi.org/10.1017/S1360223}{\textstyletext{10.1017/S1360223}}.},
	number = {1},
	pages = {67--92},
	urldate = {067432200},
	volume = {27},
	year = {2023}
}


@incollection{Sharoff2006,
	address = {Amsterdam},
	author = {Sharoff, Serge},
	booktitle = {\textit{Corpus linguistics around the world}},
	editor = {Andrew Wilson, Dawn Archer and Paul Rayson},
	pages = {167--180},
	publisher = {Rodopi},
	title = {Methods and tools for development of the {Russian} Reference Corpus},
	year = {2006}
}


Sims, Andrea D. 2023. Defectiveness. In Peter Ackema, Sabrina Bendjaballah, Eulàlia Bonet, \& Antonio Fábregas (eds) \textit{The Wiley Blackwell companion to morphology}. Wiley Blackwell. \href{https://doi.org/10.1002/9781119693604.morphcom022}{{doi: 10.1002/9781119693604.morphcom022}}.

@incollection{Sims-WilliamsSims-Williams2021,
	address = {Amsterdam},
	author = {Sims-Williams, Helen and Matthew Baerman},
	booktitle = {\textit{Lost in change: {{C}}auses and processes in the loss of grammatical elements and constructions}},
	editor = {Svenja Kranich and Tine Breban},
	note = {org/10.1075/slcs.218.},
	pages = {21–49},
	publisher = {John Benjamins doi},
	title = {A typological perspective on the loss of inflection},
	year = {2021}
}


@book{SlioussarSlioussar2015,
	address = {In},
	author = {Slioussar, N. and M. Samoilova},
	note = {\href{},
	publisher = {\textit{Proceedings of the conference ‘\citealt{Dialogue2015}’}},
	title = {Častotnosti različnyx grammatičeskix xarakteristik i okončanij u suščestvitel′nyx russkogo jazyka},
	url = {http://www.dialog-21.ru/digests/dialog2015/materials/pdf/SlioussarNASamoilovaMV.pdf}{{www.dialog-21.ru/digests/dialog2015/materials/pdf/SlioussarNASamoilovaMV.pdf}}},
	year = {2015}
}


Souag, Lameen. 2025. The diachronic stability of uninflectedness in Berber. In Enrique Palancar \& Sebastian Fedden (eds) \textit{Uninflectedness}. Berlin: Language Science Press.

@book{SpauldingSpaulding1994,
	address = {Ukarumpa via Lae},
	author = {Spaulding, Craig and Pat Spaulding},
	number = {41},
	publisher = {Summer Institute of Linguistics},
	series = {Data Papers on Papua New Guinea Languages},
	title = {\textit{Phonology and grammar of {Nankina}}},
	year = {1994}
}


@incollection{Spencer2005,
	address = {Stanford},
	author = {Spencer, Andrew},
	booktitle = {\textit{Morphology and the web of grammar: {{{E}}}ssays in memory of Steven G. {{L}}apointe} ({Stanford} Studies in Morphology and the Lexicon)},
	editor = {C. Orhan Orgun and Peter Sells},
	pages = {95--138},
	publisher = {CSLI Publications},
	title = {Towards a typology of ‘mixed categories’},
	year = {2005}
}


@incollection{Spencer2009,
	address = {Berlin},
	author = {Spencer, Andrew},
	booktitle = {\textit{On Inflection} (Trends in Linguistics, Studies and Monographs 184)},
	editor = {Patrick O. Steinkrüger and Manfred Krifka},
	pages = {173–218},
	publisher = {Mouton de Gruyter},
	title = {Realization-based morphosyntax: {{T}}he {German} genitive},
	year = {2009}
}


@incollection{Spencer2020,
	address = {Cambridge: Cambridge University Press. doi},
	author = {Spencer, Andrew},
	booktitle = {\textit{Complex Words: {{A}}dvances in Morphology}},
	editor = {Lívia Körtvélyessy and Pavol Štekauer},
	note = {1017/80643.009.},
	pages = {142–58},
	publisher = {10},
	title = {{U}ninflectedness: {{{U}}}ninflecting, uninflectable and uninflected words, or the complexity of the simplex},
	urldate = {97811087},
	year = {2020}
}


Spencer, Andrew. 2025. Some concepts and consequences of uninflectability. In Enrique Palancar \& Sebastian Fedden (eds) \textit{Uninflectedness}. Berlin: Language Science Press.

@book{SpencerSpencer2012,
	address = {Cambridge},
	author = {Spencer, Andrew and Ana R. Luís},
	publisher = {Cambridge University Press},
	sortname = {Spencer, Andrew and Ana R. Luis},
	title = {\textit{Clitics: {{A}}n introduction}},
	year = {2012}
}


@book{\textstyletext{Stefanović2019,
	address = {\textit{L}}\emph{es numéraux en bosniaque, croate, monténégrin, serbe.} Paris},
	author = {\textstyletext{Stefanović, Aleksandar},
	publisher = {Institut d’études slaves},
	sortname = {\textstyletext{Stefanovic, Aleksandar},
	year = {2019}
}


@book{Stevanović1975,
	address = {Belgrade},
	author = {Stevanović, Mihailo},
	publisher = {Naučna knjiga},
	sortname = {Stevanovic, Mihailo},
	title = {\textit{Savremeni srpskohrvatski jezik (Gramatički sistemi i književnojezička norma): {{I}}: {{{U}}}vod, fonetika, morfologija.} 3rd edition},
	year = {1975}
}


@misc{Strange1972,
	author = {Strange, David},
	note = {Ms. Summer Institute of Linguistics. \url{},
	title = {Dano noun inflection},
	url = {https://www.sil.org/resources/archives/31207}},
	year = {1972}
}


@incollection{Stump2013,
	address = {Oxford},
	author = {Stump, Gregory T.},
	booktitle = {\textit{Periphrasis:} \textit{The role of syntax and morphology in paradigms} (Proceedings of the {British} Academy 180)},
	editor = {Marina Chumakina and Greville G. Corbett},
	pages = {105--138},
	publisher = {British Academy and Oxford University Press},
	title = {Periphrasis in the {Sanskrit} verb system},
	year = {2013}
}


@incollection{\textstylengbinding{Stump2017,
	address = {Berlin: Language Science Press. DOI},
	author = {\textstylengbinding{Stump, Gregory.}},
	booktitle = {\textit{On looking into words (and beyond)}\textstylengbinding{\textit{: {{S}}tructures, Relations, Analyses}}},
	editor = {Claire Bowern, Laurence Horn and Raffaella Zanuttini},
	note = {5281/zenodo.495457.},
	pages = {421--440},
	publisher = {10},
	sortname = {\textstylengbinding{Stump, Gregory.}},
	title = {{Rules} and blocks},
	year = {2017}
}


@book{\textstyleweblinkstext{Sussex\textstyleweblinkstext{Sussex2006,
	address = {\textit{The Slavic languages,} Cambridge},
	author = {\textstyleweblinkstext{Sussex, Roland and Paul V. Cubberley},
	note = {}},
	publisher = {Cambridge University Press},
	sortname = {\textstyleweblinkstext{Sussex, Roland and Paul V. Cubberley},
	year = {2006}
}


@book{Swan2002,
	address = {Bloomington, IN},
	author = {Swan, Oscar E.},
	publisher = {Slavica},
	title = {\textit{A Grammar of contemporary {Polish}}},
	year = {2002}
}


@article{Testelets2010,
	author = {Testelets, Yakov},
	journal = {\textit{Jazyk i rečevaja dejatel′nost′}},
	note = {. St Petersburg State University, Faculty of Philology.},
	pages = {126--143},
	title = {Case deficient elements and the direct case condition in {Russian}},
	volume = {9},
	year = {2010}
}


@book{Thomas1983,
	address = {\textit{Canadian Slavonic Papers / Revue Canadienne des Slavistes} Vol. 25(1) (=Canadian Contributions to the IX International Congress of Slavists},
	author = {Thomas, George},
	publisher = {\citealt{Kiev1983}), 180-205},
	title = {A comparison of the morphological adaptation of loanwords ending in a vowel in contemporary {Czech}, {Russian}, and {Serbo-Croatian}},
	year = {1983}
}


Thornton, Anna M. 2019a. How (non-)canonical is Italian morphology? In Matthew Baerman, Oliver Bond \& Andrew Hippisley (eds) \textit{Morphological Perspectives: Papers in Honour of Greville G. Corbett}, 65-99. Edinburgh: Edinburgh University Press.

@incollection{Thornton2019b,
	address = {Cham: Springer. doi},
	author = {Thornton, Anna M.},
	booktitle = {\textit{Competition in inflection and word-formation}},
	editor = {Franz Rainer, Francesco Gardani, Wolfgang U. Dressler and Hans Christian Luschützky},
	note = {1007/978-3-030-02550- 2\_9.},
	pages = {223–258},
	publisher = {10},
	title = {Overabundance: {{A}} Canonical Typology},
	year = {2019}
}


Thornton, Anna M. \& Paolo D’Achille. 2025. Uninflectedness in Italian nouns and adjectives. In Enrique Palancar \& Sebastian Fedden (eds) \textit{Uninflectedness}. Berlin: Language Science Press.

@book{Timberlake2004,
	address = {A reference grammar of Russian}. Cambridge},
	author = {Timberlake, Alan},
	publisher = {Cambridge University Press},
	title = {\textit{},
	year = {2004}
}


@incollection{Toivonen2007,
	address = {Amsterdam: John Benjamins. \href{https://doi.org/10.1075/cilt.288.09toi}{doi},
	author = {Toivonen, Ida},
	booktitle = {\textit{Saami linguistics}},
	editor = {Ida Toivonen and Diane Nelson},
	note = {1075/cilt.288.09toi}.},
	pages = {227--258},
	publisher = {10},
	title = {Verbal agreement in Inari Saami},
	year = {2007}
}


@article{Unbegaun1947,
	author = {Unbegaun, Boris O.},
	journal = {\textit{Revue des études slaves}},
	pages = {130--145},
	title = {Les substantifs indéclinables en russe},
	volume = {23},
	year = {1947}
}


@misc{Wang2014,
	author = {Wang, Qian},
	note = {MA thesis, University of Oregon. \url{https://core.ac.uk/download/pdf/.pdf}.~},
	title = {\textit{Gender assignment of {Russian} indeclinable nouns}},
	urldate = {36692535},
	year = {2014}
}


@book{Webb2013,
	address = {London},
	author = {Webb, Jeremy (editor)},
	publisher = {Profile Cape},
	sortname = {Webb, Jeremy (editor)},
	title = {\textit{Nothing: {{F}}rom absolute zero to cosmic oblivion – amazing insights into nothingness}},
	year = {2013}
}


@book{WechslerWechsler2003,
	address = {Stanford},
	author = {Wechsler, Stephen and Larisa Zlatić},
	publisher = {CSLI},
	sortname = {Wechsler, Stephen and Larisa Zlatic},
	title = {\textit{The many faces of agreement}},
	year = {2003}
}


@incollection{Windschuttel2019,
	address = {Stanford},
	author = {Windschuttel, Glenn A.},
	booktitle = {\textit{Proceedings of the {LFG}’19 conference, {Australian} National University}},
	editor = {Miriam Butt, Tracy Holloway King and Ida Toivonen},
	note = { http://csli-publications.stanford.edu/LFG/2019.},
	pages = {334–352},
	publisher = {CSLI},
	title = {Morphology or syntax: {{T}}he two types of non-agreeing verb},
	year = {2019}
}


@article{Worth1966,
	author = {Worth, Dean S.},
	journal = {\textit{Zbornik za filologiju i lingvistiku}},
	pages = {10--16},
	title = {On the stem/ending boundary in {Slavic} indeclinables},
	volume = {9},
	year = {1966}
}


@incollection{Zaliznjak1973,
	address = {Moscow},
	author = {Zaliznjak, Andrej A.},
	booktitle = {\textit{Problemy grammatičeskogo modelirovanija}, 53–87. {{M}}oscow: {{N}}auka. [Reprinted in Andrej A. {{Z}}aliznjak \textit{Russkoe imennoe slovoizmenenie: {{S}} priloženiem izbrannyx rabot po sovremennomu russkomu jazyku i obščemu jazykoznaniju}},
	editor = {Andrej A. Zaliznjak},
	note = {]},
	pages = {613–647},
	publisher = {Jazyki slavjanskoj kul′tury, 2002},
	title = {O ponimanii termina ‘padež’ v lingvističeskix opisanijax},
	year = {1973}
}


@book{Zwicky1986,
	address = {\textit{Indiana University Linguistics Club Twentieth Anniversary} \textit{Volume}, 85\nobreakdash-106. Bloomington, IN},
	author = {Zwicky, Arnold M.},
	publisher = {Indiana University Linguistics Club},
	title = {Imposed versus inherent feature specifications, and other multiple feature markings},
	year = {1986}
}


@incollection{Zwicky1996,
	address = {Oxford},
	author = {Zwicky, Arnold M.},
	booktitle = {\textit{Concise encyclopedia of syntactic theories}},
	editor = {Keith Brown and Jim Miller},
	pages = {300--305},
	publisher = {Elsevier Science},
	title = {Syntax and phonology},
	year = {1996}
}


\end{verbatim}
\sloppy\printbibliography[heading=subbibliography,notkeyword=this]
\end{document}